\documentclass[12 pt]{report}
	\author{ }
	\usepackage[frenchb]{babel}
	\usepackage[utf8]{inputenc}  
	\usepackage[T1]{fontenc}
	\usepackage{amssymb}
	\usepackage[mathscr]{euscript}
	\usepackage{stmaryrd}
	\usepackage{amsmath}
	\usepackage{arcs}
	\usepackage{tikz}
	\usepackage[all,cmtip]{xy}
	\usepackage{amsthm}
	\usepackage{varioref}
\DeclareSymbolFont{yhlargesymbols}{OMX}{yhex}{m}{n}
\DeclareMathAccent{\arc}{\mathord}{yhlargesymbols}{"F3}
	\usepackage{geometry}
	\geometry{a4paper, head = 15pt, left = 1 cm, right = 1cm}
	\usepackage{lmodern}
	\everymath{\displaystyle}
	\usepackage{hyperref}
	\usepackage{array} 
	 \usepackage{float} 
	\theoremstyle{plain} 
	\newcounter{n}
	\setcounter{n}{1} 
	\newcommand{\cm}{\textbf{ cm}}
	\newcommand\ds\displaystyle
	\renewcommand{\it}{\item[$ \mathbf{\bullet\ \then}.$]\stepcounter{n} }
	\newenvironment{calcul}{\item[$\mathbf{\then}.$] \stepcounter{n}\hfil $\displaystyle  }{.$\hfil   }
	\newcommand\syst[4]{
 \begin{calcul} \left\{ \begin{array}{l} #1 = #2, \\ #3 = #4.\end{array}\right\end{calcul}}
	\newcommand\intercalc{\end{calcul}\begin{calcul}}
	
	\title{\Large Exercices tirés du manuel de 4e de C. Lebossé et C. Hémery, éd. Hachette\footnote{Vous pouvez me signaler les erreurs par courriel : \href{mailto:leturcq.d@orange.fr}{lien de contact.} 
	}
	}
	\date{1963}
	\begin{document}
	\maketitle 
	
	
	
\part{Arithmétique}
 \chapter{Puissances}
 
 \begin{itemize}
 \item Effectuer les puissances suivantes : \begin{itemize}
 \it $2^7$.
 \it $3^5$. 
 \it $5^4$. 
 \it $4^3$.
 \it $10^3$. 
 \it $10^7$. 
 \it $12^3$. 
 \it $17^4$. 
 \end{itemize}
 \item Effectuer les produits suivants : \begin{itemize}
 \it $2^4\times 2^3$.
 \it $10^5\times 10$. 
 \it $5^3\times 5^2$. 
 \it $2^4\times2^5\times2^3$.
 \it $7^4\times7^2\times7$.
 \it $10^5\times 10^3\times 10$. 
 \it $a^5\times a^2$.  
 \it $x^7\times x^6$. 
 \it $b^8\times b^2$. 
 \it $a^4\times a^5\times a^2$. 
 \it $x\times x^7\times x^2$. 
 \it $y^3\times y^4\times y^8$. 
 \it $(2^4)^3$.
 \it $(10^3)^3$.
 \it $(5^2)^2$. 
 \it $(a^5)^2$.
 \it $(x^7)^3$. 
 \it $(y^{25})^4$.
\end{itemize}  
\item Calculer les quotients suivants : \begin{itemize}
\it $\frac{2^{13}}{2^{10}}$.
\it $\frac{17^{41}}{17^{38}}$.
\it $\frac{10^{13}}{10^2}$. 
\it $\frac{5^7}{5^7}$. 
\it $\frac{a^{15}}{a^6}$.
\it $\frac{b^{17}}{b^{12}}$.
\it $\frac{x^5}{x^4}$. 
\it $\frac{x^3}{x^3}$.
\end{itemize}
\item Effectuer les puissances suivantes, de deux façons différentes : \begin{itemize}
\it $(2\times3\times5)^2$.
\it $(4\times 7\times 11)^3$.
\it $(2^2\times3\times5)^3$.
\it $(2^2\times7\times13)^4$.
\end{itemize}
\item Calculer les puissances suivantes : 
\begin{itemize}
\it $\left(\frac5{12}\right)^2$. 
\it $\left(\frac7{13}\right)^2$. 
\it $\left(\frac{14}{25}\right)^2$. 
\it $\left(\frac47\right)^3$. 
\it $\left(\frac58\right)^3$. 
\it $\left(\frac{11}{20}\right)^3$.
\it $(0,5)^4$. 
\it $(1,25)^3$. 
\it $(0,75)^5$.
\end{itemize}
\item Effectuer les produits suivants : 
\begin{itemize}
\it $(2\times 3^2\times 7)\times (2^3\times 3 \times 5)$.
\it $(5^2\times 7\times 13)\times(5\times 7\times 11^2)$. 
\it $(2^3\times 3^2\times 5) \times ( 2^2\times 7\times 11)$. 
\it $(2\times 5 \times 7)^2\times (3\times 4)^2$.
\end{itemize}
\item Simplifier les quotients exacts suivants : 
\begin{itemize}
\it $\frac{2^5\times7^2\times17}{2^3\times7^2\times17^2}$.
\it $\frac{3^5\times 4^7\times 13}{3^6\times 4^6\times 11}$.
\it $\frac{11^2\times 13\times 29^3}{11\times 13^2\times 29^3}$. 
\it $\frac{19^{13}\times 31^{14} \times 5}{19^{12}\times 31^{15}\times2}$. 
\end{itemize}
 \end{itemize}
 \chapter{Nombres premiers}
 \begin{itemize}
 \it Établir la liste et le nombre des diviseurs de $54$. Grouper par deux les diviseurs dont le produit est $54$ et montrer qu'il suffit de rechercher le plus petit nombre de chaque groupe. \end{itemize}
 \begin{itemize}
 \item Reprendre le même problème pour les nombres :  \begin{itemize}
 \it $80$.
 \it $108$.
 \it $128$. 
 \it $252$. 
 \it $84$. 
 \it $250$.
 \it $288$. 
 \it $315$. 
 \it $36$. 
 \it $100$.
 \it $144$. 
 \it $225$.
 \end{itemize}
 \item Reconnaître si les nombres entiers suivants sont premiers et donner s'il y a lieu leur plus petit diviseur premier : \begin{itemize}
 \it $79$. 
 \it $107$. 
 \it $143$. 
 \it $173$. 
 \it $83$. 
 \it $149$. 
 \it $181$. 
 \it $221$. 
 \it $89$. 
 \it $167$. 
 \it $187$. 
 \it $241$. 
 \it $97$. 
 \it $179$. 
 \it $193$. 
 \it $283$. 
 \end{itemize}
 \it Montrer que tout nombre premier supérieur à $5$ est obligatoirement terminé par $1$, $3$, $7$ ou $9$. 
 \item Décomposer en facteurs premiers les nombres entiers suivants : 
 \begin{itemize}
 \it $108$. 
 \it $144$. 
 \it $2~520$. 
 \it $8~000$. 
 \it $84$. 
 \it $250$. 
 \it $864$. 
 \it $5~740$. 
 \it $176$.
 \it $294$. 
 \it $7~920$. 
 \it $1~053$. 
 \it $36\times 42$. 
 \it $72\times 77$. 
 \it $108\times 75$. 
 \it $84\times 25\times 121$.
 \it $36\times 27\times 143$. 
 \it $65\times 49\times 24$. 
 \it $108^2$. 
 \it $252^3$. 
 \it $24^2\times 33^3 \times 11^5$. 
 \end{itemize}
 \item Calculer les nombres : 
 \begin{itemize}
 \it $2^2\times 3\times 7$.
 \it $2\times 3^2\times 5\times 7$.
 \it $2^3\times 3\times 5\times 11$.
 \it $3^2\times 5^2\times 11$. 
 \it $3^2\times 7\times 11\times 13$. 
 \it $2^2\times 7^3\times 11\times 17$.
 \end{itemize}
 \item Effectuer en laissant les résultats sous forme
 décomposée : 
 \begin{itemize}
 \it $(2^2\times3^4\times 5)\times(2\times3\times 7^2)$. 
 \it $(2^2\times 3^4\times 5) \times (3^2\times 7
 \times11^3)\times(5\times 11^2)$. 
 \it $(2^4\times 3^2\times 7\times 11^3)^2$. 
 \it $(2\times 3^4\times 7^3\times 11^2)^3$. 
 \it $(2^5\times 3^2\times 5^2\times 7^2)\div (2^3\times 5\times 7^2)$.
 \it $(2^7\times 3^2\times 5^4\times 11) \div (2^7\times 3\times 5^2)$.
 \end{itemize}
 \end{itemize}
  
 \chapter{Plus grand commun diviseur - plus petit commun multiple} 
 
 \begin{itemize}
 \item Calculer le plus grand commun diviseur (PGCD) des nombres suivants : 
 \begin{itemize}
 \it $168$ et $360$. 
 \it $252$ et $684$. 
 \it $336$ et $462$. 
 \it $1~840$ et $1~260$. 
 \it $18~150$ et $23~850$. 
 \it $33~390$ et $58~800$. 
 \it $315$, $819$ et $924$. 
 \it $252$, $693$ et $945$. 
 \it $2~520$, $3~150$ et $4~410$. 
 \it $7~560$, $10~080$ et $12~096$. 
 \end{itemize}
 \item Établir la liste des diviseurs communs aux nombres : 
 \begin{itemize}
 \it $4~200$ et $5~880$. 
 \it $1~440$ et $1~764$. 
 \it $3~780$, $4~320$ et $5~184$. 
 \it $10~584$, $11~520$ et $13~104$. 
 \end{itemize}
 \item Calculer le plus petit commun multiple (PPCM) et
 les trois multiples communs les plus simples de : 
 \begin{itemize}
 \it $360$ et $504$. 
 \it $252$ et $672$. 
 \it $972$ et $1~134$. 
 \it $720$ et $900$. 
 \it $168$, $252$ et $336$. 
 \it $120$, $180$ et $270$. 
 \end{itemize}
 \it \begin{enumerate}
 \item Calculer le PGCD et le PPCM des nombres $576$ et $1~080$. 
 \item Comparer le produit des deux résultats au produit des deux nombres. Énoncer le résultat obtenu.
 \end{enumerate}
 \it \begin{enumerate}
 \item Calcule le PGCD et le PPCM des nombres $99$ et $140$. 
 \item Quel est le PPCM de deux nombres premiers entre eux ? 
 \end{enumerate}
 \it \begin{enumerate}
 \item Calculer le PGCD et le PPCM des nombres $144$ et 
 $180$. 
 \item Que deviennent les résultats précédents lorsqu'on
 multiplie (ou divise) les deux nombres par $6$ ?
 Généraliser. 
 \end{enumerate}
 \it Démontrer que, si un nombre en divise deux autres, il divise leur somme, leur différence, et le reste de 
 leur division. 
 \it La division de deux nombres se fait exactement. Quel
 est leur PGCD et quel est leur PPCM ? Exemple : $6~375$
 et $375$. 
 \it Montrer que l'ensemble des diviseurs communs à deux nombres est le même que celui du plus petit de ces nombres et du reste de leur division. Que peut-on dire des PGCD ? \\ 
 \emph{Application.} — Remplacer la recherche du PGCD de $792$ et $240$ par la recherche du PGCD de deux nombres plus simples. Répéter cette opération afin d'obtenir un
 PGCD évident. (\emph{Méthode des divisions successives}.)
 \it Utiliser la méthode indiquée au numéro précédent pour la recherche du PGCD des nombres $2~021$ et $2~679$. 
 \it Trouver deux nombres non divisibles l'un par l'autre sachant que leur PGCD est égal à $336$ et leur somme 
 égale à $2~688$. 
 \it Par quel nombre inférieur à $100$ faut-il diviser $29~687$ et $35~312$ pour obtenir pour restes respectifs $47$ et $32$ ? 
 Quels sont les quotients ? 
 \it En divisant $809$ et $1~024$ chacun par un certain nombre, on trouve le même quotient et pour restes respectifs $27$ et $35$. Reconstituer les deux divisions. 
 \it On a planté des arbres également espacés sur le pourtour d'un terrain triangulaire dont les côtés mesurent $144$ m, $180$ m, et $240$ m. Sachant qu'il y a un arbre à chaque sommet et que la distance de deux arbres consécutifs est comprise entre $4$ mètres et $10$ mètres, calculer le nombres d'arbres plantés.
 \it Un ouvrier a touché pour trois mois successifs : $462$ F, $528$ F, et $594$ F. Trouver son salaire journalier sachant que c'est un nombre entier de francs 
 compris entre $20$ F et $30$ F. Trouver le nombre
 de jours de travail effectués chaque mois.
 \it On a fait carreler une pièce rectangulaire de $4,2$ m sur $2,24$ m. Sachant que les carreaux employés ont un côté compris entre $10$ cm et $25$ cm, calculer la longueur de leur côté et leur nombre. 
 \it On veut partager en coupons d'égale longueur quatre pièces d'étoffe mesurant respectivement $17,50$ m, $28$ m, $31,50$ m, et $42$ m. Trouver la plus grande longueur possible pour chaque coupon et le nombre total de ces coupons. 
 \it Deux règles égales de $504$ mm de longueur sont graduées l'une en $72$ parties et l'autre en $126$ parties. On fait coïncider leurs extrémités. Déterminer les traits de graduation qui coïncident. 
 \it Trouver les multiples communs à $6$, $8$ et $10$ compris entre $500$ et $1~000$. 
 \it Trouver les trois nombres les plus simples divisibles par les $10$ premiers nombres entiers. 
 \it Quel est le plus petit nombre qui donne $7$ pour reste quand on le divise par $12$, par $15$ ou par $16$ ? 
 \it Trouver un nombre qui donne $16$ pour reste quand on
 le divise par $24$ ou par $32$, et $8$ pour reste lorsqu'on le divise par $20$. 
 \it Trouver le plus petit nombre qui, divisé par $5$, $6$ ou $8$, donne respectivement pour restes $4$, $5$ et $7$. 
\it Deux cyclistes roulent dans le même sens sur une piste. Le premier fait un tour en $1$ minute $45$ seconde et le second en $1$ minute $36$ secondes. Sachant qu'ils sont partis ensemble de la ligne de départ, on demande après combien de temps passeront-ils ensemble cette ligne de départ. Combien de tours chacun d'eux aura-t-il effectués ? 
\it Des pavés rectangulaires qui ont $12$ cm de large et
$21$ cm de long ont servi à paver entièrement une place carrée. Calculer le côté de cette place sachant qu'il mesure un nombre entier de mètres compris entre $30$ mètres et $60$ mètres. 
\it Trois règles graduées de $960$ mm de long sont placées côté à côte de façon que leurs extrémités coïncident. Les divisions ont pour longueurs respectives $10$ mm, $12$ mm et $16$ mm. Déterminer les traits de divisions qui coïncident sur les trois règles. 
\it Une personne a acheté un certain nombre entier de mètres d'étoffe à $6,50$ F le mètre. Elle paye exactement avec des billets de $10$ F. Sachant que la dépense totale n'excède pas $200$ F, trouver le nombre de mètres achetés.
\it Un enfant compte ses timbres-poste par $12$, par $16$ et par $20$. Il lui en reste $8$ à chaque fois. En les comptant par $13$, il n'en reste plus. Combien possède-t-il de timbres ? 
\it En comptant les élèves d'une école par $9$, $10$ ou $12$, il en reste respectivement $8$, $9$ et $11$. En les comptant par $11$, il n'en reste pas. Trouver le nombre
d'élèves de l'école. 
\it Des autobus partent d'un même point dans quatre directions différentes. Les départs se font respectivement dans chaque direction toutes les $5$ minutes, $8$ minutes, $12$ minutes et $18$ minutes. Un départ simultané a lieu à $7$ heures le matin. Quelles sont les heures des autres départs simultanés de la journée ? 

 \end{itemize}
 
 \chapter{Application aux fractions} 
 \begin{itemize}
 \item Simplifier les fractions suivantes : 
 \begin{itemize}
 \it $\frac{1~815}{2~385}$ et $\frac{2~184}{7~560}$.
 \it $\frac{8~613}{9~009}$ et $\frac{12~012}{25~025}$.
 \it $\frac{3~339}{5~880}$ et $\frac{17~226}{18~810}$.
 \it $\frac{6~720}{9~405}$ et $\frac{15~120}{30~576}$.
 \it $\frac{5~292}{8~400}$ et $\frac{9~360}{18~144}$. 
 \it $\frac{5~203}{6~149}$ et $\frac{43~659}{68~607}$.
 \it $\frac{168\times462}{336\times360}$ et $\frac{252\times684}{840\times1~260}$. 
 \it $\frac{507\times 451}{861\times 396}$ et $\frac{315\times 693}{924\times 504}$. 
 \it $\frac{441\times 1~815\times 3~339}{588\times 3~150\times 2~385}$. 
 \it $\frac{756\times 336\times 2~205}{252\times 924\times12~096}$. 
 \end{itemize}
 \item Effectuer les opérations suivantes : 
 \begin{itemize}
 \it $\frac{90}{189}+\frac{45}{84}-\frac{75}{126}$.
 \it $\frac{51}{56}+\frac{88}{231}-\frac{26}{39}$.
 \it $\frac{77}{35}-\frac{80}{66}-\frac{56}{105}$. 
 \it $\frac{91}{52}+\frac9{84}-\frac{55}{105}$.
 \it $\frac{96}{160}+\frac{57}{105}-\frac{39}{77}$. 
 \it $\frac{171}{76} -\frac{161}{147} - \frac{35}{60}$.
 \it $\left(\frac{168}{192}-\frac{50}{112}+\frac{65}{273}\right)\times \frac34$.
 \it $\left(\frac{85}{153}+\frac{21}{60}-\frac{35}{225}\right)\times \frac25$. 
 \it $\left(\frac{18}{72}-\frac{23}{115}+\frac{45}{63}\right)\div \frac27$.
 \it $\left(\frac{100}{275}-\frac{45}{189}+\frac{60}{198}\right)\div\frac3{11}$. 
 \end{itemize}
 \it Trouver une fraction égale à $\frac{52}{117}$ et dont la somme des termes soit égale à $325$.
 \it Trouver une fraction égale à $\frac{44}{121}$ dont la différence des termes soit égale à $357$. 
 \it La somme de deux fractions dont l'une est les $\frac37$ de l'autre est égale à $\frac{40}{21}$. Calculer
 les deux fractions.
 \it Trouver une fraction égale à $\frac{192}{300}$ dont le PGCD des termes soit égal à $13$. 
 \it Trouver une fraction égale à $\frac{132}{385}$ dont
 le PPCM des termes soit $2~100$. 
 \it Trouver deux fractions sachant que leur quotient est égal à $\frac{23}5$ et que leur somme est égale à $\frac{35}3$.
 \it Trouver deux fractions sachant que leur quotient
 est égal à $\frac{17}8$ et que leur différence est égale à $\frac{23}{28}$. 
 \it Trouver les fractions égales à $\frac{126}{192}$ et :\begin{enumerate}
 \item Dont les termes soient inférieur à ceux de la fraction proposée ; 
 \item Dont le numérateur soit inférieur à $400$ et le dénominateur supérieur à $500$ ; 
 \item Dont la différence des termes soit égale à $132$.
 \end{enumerate}
 \it Trouver deux fractions ayant pour numérateur $1$, dont les dénominateurs soient deux nombres entiers consécutifs et comprenant entre elles la fraction $\frac{13}{84}$. 
 \it Deux règles graduées placées côte à côte de façon que leurs origines coïncident mesurent respectivement $1,40$ m et $1,68$ m. La première est partagée en $48$ parties égales et la seconde en $32$ parties égales. Déterminer par leurs numéros les traits de deux graduations qui coïncident. 
 \it Deux pièces d'étoffe mesurent respectivement $\frac{225}8$ de mètre et $\frac{175}{12}$ de mètre. 
 On veut les découper en coupons d'égale longueur. quelle est la plus grande longueur possible pour chacun de ces coupons et combien de coupons chacune de ces pièces fournira-t-elle ? 
 \it Trouver la plus petite fraction dont les quotients par $\frac{28}{45}$ et par $\frac{40}{63}$ soient des nombres entiers. Comparer les termes de la fraction irréductible trouvée au PPCM des numérateurs et au PGCD des dénominateurs des deux fractions proposées. 
 \it Soit les fractions $\frac{63}{22}$ et $\frac{99}{20}$. Trouver la plus grande fraction qui soit contenue un nombre entier de fois dans chacun de ces deux fractions. Comparer la fraction trouvée au PGCD des numérateurs et au PPCM des dénominateurs des deux fractions proposées.
 \it Réduire la fraction $\frac{168}{315}$ à sa plus simple expression. \\ Trouver ensuite toutes les fractions égales à la fraction donnée, et à termes plus petit. Quel est le nombre de ces fractions ?\\
 Déterminer une fraction égale à $\frac{168}{315}$ et dont le numérateur et le dénominateur ont pour somme $8~303$. 
 \it Calculer l'excès de l'unité sur la fraction $\frac8{11}$. \\ On ajoute $5$ à chacun des deux termes de cette fraction (on remarquera que la différence des deux termes ne change pas) : calculer l'excès de l'unité sur la fraction ainsi obtenue. Dire d'après cela si la fraction $\frac8{11}$ augmente ou diminue lorsqu'on ajoute un même 
 nombre à ses deux termes. \\En employant un procédé analogue, dire si la fraction $\frac{15}7$ augmente ou 
 diminue lorsqu'on ajoute un même nombre à ses deux termes.
 \it Un marchand revend à raison de $9$ F le mètre une pièce d'étoffe qu'il a achetée à un prix inconnu. \\
 Il vend une première fois les $\frac38$ de la pièce, une deuxième fois les $\frac35$ du reste et une troisième fois la moitié du nouveau reste. Ces trois ventes ont déjà produit une somme égale au prix d'achat total de la pièce, augmenté de $9$ F. Dans une quatrième vente, le marchand vend le reste de la pièce, et son bénéfice total est $144$ F. \\Calculer la longueur de la pièce et le prix d'achat du mètre. 
 \it Une pièce de tissu de $84$ mètres a été vendue à trois acheteurs. Le premier a eu les $\frac25$ de la pièce ; le second a eu une part égale aux $\frac57$ de la
 part vendue au premier ; le troisième a eu le reste.
 Calculer la longueur du tissu vendu à chaque acheteur.\\
 Le premier a payé $45$ F par mètre, les deux autres ont payé $70$ F par mètre. 
Calculer le prix d'achat de la pièce, sachant que le bénéfice du marchand est le $\frac13$ du prix d'achat.
 \end{itemize}
 \part{Algèbre}
 
 \chapter{Nombres relatifs}
  \begin{itemize}
  \it Peut-on mesurer à l'aide des nombres relatifs les 
  grandeurs suivantes : fortune d'une personne, gains d'un
  joueur, vitesse d'un mobile sur un axe, date d'un évènement ?
  \it Placer sur un axe $(xy)$ les points suivants donnés par leurs abscisses, exprimées en centimètres : $A(+3); B(-2); C(-5); D(+4)$.\\En déduire les valeurs de $\overline{AB}$, $\overline{AC}$, $\overline{AD}$, $\overline{BC}$, $\overline{BD}$ et $\overline{CD}$.
  \it Sur un axe $(x'x)$ d'origine $O$, on construit les points $A(+3)$, $B(+11)$, et $M(+7)$. Calculer les mesures algébriques $\overline{AB}$, $\overline{AM}$, $\overline{MB}$ et $\overline{OM}$. Que représente le point $M$ pour le segment $[AB]$ et que peut-on dire des vecteurs $\overrightarrow{MA}$ et $\overrightarrow{MB}$ ? 
  \it On construit sur un axe $(x'Ox)$ les points 
  $A(+5), B(+9), C(-1)$ et $D(-5)$. Comparer les vecteurs $\overrightarrow{AB}$ et $\overrightarrow{DC}$ ainsi que les vecteurs $\overrightarrow{AD}$ et $\overrightarrow{BC}$. Que représente le point $M(+2)$ pour chacun des segments $[AC]$ et $[BD]$ ? 
  \it On place sur un axe $(xy)$ deux points $A$ et $B$ d'abscisses $+1$ et $-3$. Quelles sont les abscisses des points qui partagent le segment $[AB]$ en huit parties égales ? 
  \it Soient $A$ et $B$ deux points d'abscisses $-2$ et $-5$. Quelle est l'abscisse du milieu $M$ de $[AB]$ ? Quelle est l'abscisse du point $N$ tel que $\overline{BN} = 2\overline{AN}$ ?
  \it Soit un point d'abscisse $+5$ ; quelle est l'abscisse du point $A'$ symétrique de $A$ par rapport à l'origine $O$ des abscisses ? Soit de même $B$ d'abscisse $-1$ et $B'$ son symétrique par rapport à $O$. Évaluer $\overline{AB}$ et $\overline{A'B'}$. 
  \it Dans l'échelle d'un thermomètre Fahrenheit la division $32$ correspond au $0^o$ centigrade et la division $212$ à $100^o$ centigrades. 
  \begin{enumerate}
  \item À combien de degrés centigrades correspond un degré Fahrenheit ? 
  \item Trouver les températures centigrades correspondant à $+50$ degrés Fahrenheit et à $+5$ degrés Fahrenheit. 
  \end{enumerate}
  \it \begin{enumerate}
  \item À combien de degrés Fahrenheit correspond un degré centigrade ?
  \item Trouver la température Fahrenheit correspondant à
  $+15^o$ centigrades et à $-15^o$ centigrades.
  \end{enumerate}
  \end{itemize}
 \chapter{Addition des nombres relatifs}

\begin{itemize}
\item Effectuer les additions suivantes : 
\begin{itemize}
\it $(+3)+(-5)+(+7)$.
\it $(+13)+(+17)+(+15)$. 
\it $(-17)+(+14)+(-41)$. 
\it $(+3)+(-3)+(+5)+(-5)$. 
\it $(-25)+(-30)+(-40)$.
\it $(-100)+(+75)+(+25)$. 
\end{itemize}
\item Effectuer les additions suivantes : 
\begin{itemize}
\it $\left(-\frac14\right)+\left(-\frac32\right)+\left(+\frac52\right)$.
\it $\left(-\frac56\right)+\left(+\frac49\right)+\left(-\frac73\right)$.
\it $(-0,75) + (+4) + \left(-\frac45\right)$.
\it $\left(+\frac47\right)+\left(-\frac{13}{21}\right)+\left(+\frac56\right)$.
\it $(-12)+(-11,57)+(+9,87)$. 
\it $ (-4,51)+\left(-\frac45\right)+\left(+\frac7{20}\right)$.
\end{itemize}
\item Effectuer les additions suivantes : 
\begin{itemize}
\it $-13+15-7-9+1$.
\it $+\frac34+\frac73+\frac5{12}+1$.
\it $-\frac12+0,4-0,7+3$.
\it $+\frac12-\frac49+\frac5{18}-\frac74$.
\it $+7-9-11+0,75$. 
\it $-4-\frac73+\frac78$.
\end{itemize}
\item Effectuer les additions suivantes : 
\begin{itemize}
\it $(+7-8+4)+(-9-13-1)+(+4-5-7)$.
\it $[(-3)+(-7)+(-1)] + [(-5)+(-9)+(+4)]$.
\it $\left(\frac{13}{56}-\frac3{14}+\frac5{28}\right)+\left(\frac47-3\right)$. 
\it $[(+5-0,7)+(+4-0,3)]+[(-13-4,7)+(-3)]$.
\it Placer sur un axe les point suivants dont les abscisses sont données en centimètres :\\ $A(-3); B(-1); C(+1); D(+2)$ et $E(+5)$. Vérifier les égalités :
\[\begin{array}{ll}
\overline{AE} = \overline{AB}+\overline{BE} & \overline{AE} = \overline{AD}+\overline{DE}\\
\overline{AE} = \overline{AB}+\overline{BC}+\overline{CE}&\overline{AE} = \overline{AD}+\overline{DB}+\overline{BE}\\
\overline{AE}=\overline{AB}+\overline{BC}+\overline{CD}+\overline{DE} & \overline{AE} = \overline{AD}+\overline{DB}+\overline{BC}+\overline{CE}
\end{array}\]
\end{itemize}
\end{itemize}
 \chapter{Soustraction des nombres relatifs}


\begin{itemize}
\item Calculer les différences : 
\begin{itemize}
\it $(-15)-(-13)$.
\it $(+4,5)-(+5,49)$. 
\it $\left(+\frac34\right)-\left(+\frac73\right).$
\it $(+12)-(+9)$.
\it $(-3,9)-(-7,4)$.
\it $\left(-\frac34\right)-\left(-\frac79\right).$
\it $(+5)-(+17)$.
\it $(-13,75)-(-4,01)$. 
\it $\left(-\frac25\right)-\left(-\frac1{13}\right)$.
\end{itemize}
\item Calculer la valeur des sommes algébriques suivantes : 
\begin{itemize}
\it $(-10)+(-15)-(-30)+(+7)-(+12)$.
\it $(-7)-(-9)-(+17)+(-41)+(+1)$.
\it $(+49)-(-63)+(-3)-(+4)$.
\end{itemize}
\item Calculer la valeur des sommes algébriques suivantes : 
\begin{itemize}
\it $(-0,7)-(-0,9)+(-13,5)-(+4)$
\it $\left(-\frac12\right)+\left(-\frac34\right) - \left(+\frac25\right)-(-1)$.
\it $\left(-\frac5{11}\right)+\left(-\frac7{22}\right)
-\left(-\frac4{33}\right) + (-1) -(+0,5)$.
\end{itemize}
\item Calculer la valeur des sommes algébriques suivantes : 
\begin{itemize}
\it $7-10+14-13-21+3$.
\it $0,75+4,7-0,7-0,5-1$.
\it $\frac23-\frac34+\frac5{12}-\frac56 - 3$. 
\end{itemize}
\item Effectuer les opérations suivantes : 
\begin{itemize}
\it $(5-3+7-1)+(-9+4-1)-(-3-7+2)$.
\it $(-14+7-3)+(5-7-8)-(-5+4+10)$.
\it $\left(-\frac23-\frac35+1\right) - \left(-\frac13+\frac45-2\right) + \left(-1+\frac73\right)$.
\end{itemize}
\item Effectuer les opérations suivantes : 
\begin{itemize}
\it $[(5-9)+(3-5)] - [(7+3-5) - (7-10)]$.
\it $[12-(14-5+0,75)] + [-15 + (3,25-2)]$.
\it $\left[1 - \left(\frac25-\frac43\right)\right] - 
\left[ \left(\frac45-1\right) - \left(\frac75+2\right)\right]$. 
\end{itemize}
\item Supprimer les parenthèses et les crochets dans les expressions suivantes : 
\begin{itemize}
\it $(a-b+c) - (d- e - f) + (b-a)$. 
\it $[(a-b)-(a-5)] + [b-7-(a-3)]$. 
\it $[12-(a-b)+6] - [15+ (b-a-13)]$. 
\end{itemize}
\it Quel est l'accroissement de la température indiquée 
par un thermomètre qui passe de $+7^o$ à $+13^o$, ou de $-3^o$ à $+1^o$, ou de $-5^o$ à $-7^o$, ou de $+3^o$ à $-1^o$ ?
\item Quels sont les intervalles de temps qui séparent les dates suivantes ? 
\begin{itemize}
\it $+2$h$15$min et $+11$h$45$ min. 
\it $-2$h$12$min et $-1$h$17$ min. 
\it $+3$h$51$min et $+7$h$17$ min. 
\it $-3$h$10$min et $-1$h$5$ min. 
\it $-7$h$14$min et $+8$h$16$ min. 
\it $-13$h$45$min et $+5$h$51$ min. 
\it $-4$h$17$min$51$sec et $+12$h$17$ min$47$sec. 
\end{itemize}
\end{itemize}
 \chapter{Produit des nombres relatifs}

\begin{itemize}
\item Effectuer : 
\begin{itemize}
\it $(+17)(-13)$.
\it $\left(+\frac23\right)\left(+\frac57\right)$.
\it $(-0,7)\left(+\frac7{11}\right)$.
\it $(-14)(-11)$.
\it $\left(-\frac{11}{11}\right)\left(-\frac34\right)$.
\it $\left(-\frac32\right)(-0,5)$.
\it $(+12)(+4)$.
\it $\left(+\frac59\right)\left(-\frac75\right)$.
\it $(+0,7)\left(-\frac49\right)$.
\end{itemize}
\item Effectuer 
\begin{itemize}
\it $(-3)(+7)(-5)(-4)$.
\it $(+1,7)(-2,5)(-3)(-1,9)$.
\it $(-1)(+1)(-1)(+1)(-1)(-1)$.
\it $\left(-\frac35\right)(+3)\left(-\frac53\right)$.
\it $\left(+\frac12\right)\left(-\frac14\right)\left(-\frac48\right)$.
\it $\left(+\frac12\right)\left(-\frac6{11}\right)\left(-\frac78\right)\left(-\frac{11}{15}\right)$.
\end{itemize}
\item Calculer de deux façons différentes les produits suivants : 
\begin{itemize}
\it $(-5+7-9-1)(-3)$.
\it $(+12-17-9)(+7)$.
\it $(-12+0,7)(+0,5)$.
\it $(-13+11,7+12,5)(-4,7)$.
\it $\left(-\frac12+\frac23+\frac35\right)(-4)$.
\it $\left(-\frac35+\frac79+\frac{11}{13}\right)\left(+\frac23\right)$.
\end{itemize}
\item Calculer de deux façons différentes les produits suivants : 
\begin{itemize}
\it $(-5+12-7)(5-7)$.
\it $\left(\frac35-\frac4{15}\right)\left(\frac27-\frac3{14}\right)$.
\it $(+5+14-13)(-7-2)$.
\it $\left(\frac59-\frac{7}{27}\right)\left(1-\frac35\right)$.
\it $(-2,5+2,7+4,9)(0,75-1)$.
\it $\left(\frac{11}{12}-1\right)\left(\frac{11}{12}+1\right)$.
\end{itemize}
\item Effectuer le plus rapidement possible les opérations suivantes : 
\begin{itemize}
\it $(5-3+2)(12-9)+(15-17)(7-10)$.
\it $(7-17)(10-12) - (15-17)(-14+11)$.
\it $(2,5+0,7)(7,9-8)-0,5(11-8)$.
\it $\left(\frac23+\frac35-1\right)\left(1-\frac79\right)+\frac45\left(17-\frac13\right)$.
\it $[(5-11)-(17-14)][12-(14-11)] + 13(17-12)$.
\it $[14-(0,7-1)]\left[0,5+\left(\frac12-3\right)\right] - [12-(2,4+4,1)][(1,7+2,3)-1]$. 
\end{itemize}
\it La vitesse d'un point mobile sur un axe est $+4$ cm par seconde, et ce mobile passe au point $O$, origine des abscisses au temps $t=0$. Quelles sont les abscisses de ce point aux temps $+3$ et $+7$ ? Quelle est la valeur algébrique du vecteur joignant les $2$ positions de ce point aux temps considérés ? Même question si la vitesse est $-4$ cm par seconde. 
\it Au temps $+5$, l'abscisse d'un mobile est $-16$ cm, au temps $-3$ elle est $+8$. Quelle est la vitesse de ce mobile sachant qu'il est animé d'un mouvement uniforme ? 
\end{itemize}
 \chapter{Division des nombres relatifs}

\begin{itemize}
\item Effectuer les divisions suivantes :
\begin{itemize}
\it $(+56):(+7)$.
\it $(-56):(-8)$.
\it $(+64):(-16)$.
\it $(-360):(+12)$.
\it $(+105):(-7)$. 
\it $(-286):(-13)$.
\it $(+221):(-17)$.
\it $(-286):(-11)$.
\it $(+96):(+8)$. 
\end{itemize} 
\item Effectuer les divisions suivantes : 
\begin{itemize}
\it $\displaystyle(+15):\left(-\frac23\right)$.
\it $\displaystyle(-7,5):\left(-\frac45\right)$.
\it $\displaystyle(-1):\left(+\frac34\right)$.
\it $\displaystyle\left(-\frac45\right):(+1,4)$.
\it $\displaystyle\left(+\frac34\right):\left(-\frac{15}8\right)$.
\it $\displaystyle\left(-\frac73\right):\left(+\frac53\right)$.
\it $\displaystyle(+1):\left(+\frac23\right)$.
\it $\displaystyle\left(-\frac{8}{25}\right):\left(-\frac{14}5\right)$.
\it $\displaystyle\left(+\frac13\right):\left(-\frac14\right)$.
\end{itemize}
\item Effectuer les divisions suivantes : 
\begin{itemize}
\it $\displaystyle\frac{-12+4-72}{4}$.
\it $\displaystyle\frac{+35-75+25}{5}$.
\it $\displaystyle\frac{-12+144-108}{-12}$.
\it $\displaystyle\frac{7-0,5+2}{-2}$. 
\it $\displaystyle\frac{9+13-15}{11}$.
\it $\displaystyle\frac{-11+15-7,8}5$. 
\end{itemize}
\item Simplifier les rapports suivants et trouver leur valeur : 
\begin{itemize}
\it $\displaystyle\frac{-63}{70}$.
\it $\displaystyle\frac{-308}{484}$.
\it $\displaystyle\frac{273}{468}$. 
\it $\displaystyle\frac{48}{-96}$.
\it $\displaystyle\frac{-169}{-26}$.
\it $\displaystyle\frac{-365}{-146}$.
\end{itemize}
\item Effectuer les opérations suivantes : 
\begin{itemize}
\it $\displaystyle\frac{-3+\frac34-\frac13}{5+\frac25-\frac23}$.
\it $\displaystyle\frac{7-\frac32+\frac38}{2-\frac52-\frac34}$.
\it $\displaystyle\frac{1-\frac57-\frac49}{1-\frac53-\frac7{15}}$.
\it $\displaystyle\frac{-3-0,5-0,75}{-\frac23-\frac34}$. 
\end{itemize}
\item Effectuer les opérations suivantes : 
\begin{itemize}
\it $\displaystyle\frac{-3+5-6}{2+7-1}\times\frac{4-3+1}{2-3-1}\times \frac{-5-7}{10}$.
\it $\displaystyle\frac{\frac12-\frac23}{\frac34-1}\times \frac{\frac56+\frac13}{1-\frac45}\times \frac{-\frac78+1}{\frac32}.$
\end{itemize}
\item Effectuer les opérations suivantes : 
\begin{itemize}
\it $\displaystyle\frac{0,5}{\frac35-2}+\frac{1}{4-\frac{11}7}-\frac3{-7+\frac13}$.
\it $\displaystyle 3- \frac{1-\frac12}{1+\frac12}$.
\it $\displaystyle \frac{\frac35}{\frac23-\frac45}-
\frac{\frac12}{\frac18-1}+\frac{\frac12}{\frac1{16}-2}$.
\it $\displaystyle\frac{1-\frac13+\frac1{1+\frac13}}{1-\frac13-\frac1{1+\frac13}}$
\end{itemize}
\item Le nombre $x$ prend toutes les valeurs entières de $-5$ à $+10$. Dresser le tableau de correspondance entre $x$ et $y$ dans les cas suivants : 
\begin{itemize}
\it $\displaystyle y=\frac{3x}5+1$.
\it $\displaystyle y=\frac{-2x}3+4$.
\it $\displaystyle y=\frac{-3x}2-4$.
\it $\displaystyle y=\frac{-7x}5-10$.
\it $\displaystyle y=\frac{2x+3}5$.
\it $\displaystyle y=\frac{-3x+1}4$.
\it $\displaystyle y=\frac{3x-7}2$.
\it $\displaystyle y=\frac{-5x-3}3$.
\end{itemize}
\end{itemize}
 \chapter{Relation de Chasles}

\begin{itemize}
\item Vérifier la relation $\overline{AC} = 
\overline{AB}+\overline{BC}$, les points $A$, $B$ et $C$ appartenant à un axe $(xy)$, dans les cas suivants : 
\begin{itemize}
\it Les abscisses des points $A$, $B$, $C$ sont $0$,
$+3$, $+5$.
\it Les abscisses des points $A$, $B$, $C$ sont $0$,
$-2$, $-4$.
\it Les abscisses des points $A$, $B$, $C$ sont $0$,
$+7$, $+2$.
\it Les abscisses des points $A$, $B$, $C$ sont $0$,
$-5$, $-3$.
\it Les abscisses des points $A$, $B$, $C$ sont $0$,
$+2$, $-3$.
\it Les abscisses des points $A$, $B$, $C$ sont $0$,
$-3$, $+2$.
\end{itemize}
\item Soit un axe $(xy)$, $O$ l'origine des abscisses, vérifier la relation $\overline{AB} = 
\overline{OB}-\overline{OA}$ et déterminer la longueur $AB$ dans les cas suivants : \begin{itemize}
\it Les abscisses de $A$ et $B$ sont $+3$, $+5$.
\it Les abscisses de $A$ et $B$ sont $-2$, $-4$.
\it Les abscisses de $A$ et $B$ sont $+7$, $+2$.
\it Les abscisses de $A$ et $B$ sont $-5$, $-3$.
\it Les abscisses de $A$ et $B$ sont $+2$, $-3$.
\it Les abscisses de $A$ et $B$ sont $-3$, $+2$.
\end{itemize}
\item Soient $A$ et $B$ deux points d'un axe, $M$ le milieu de $[AB]$; vérifier la relation $\displaystyle\overline{OM} = \frac{\overline{OA}+\overline{OB}}2$
dans les cas suivants : \begin{itemize}
\it Les abscisses de $A$ et $B$ sont $+5$, $+9$.
\it Les abscisses de $A$ et $B$ sont $-5$, $+9$.
\it Les abscisses de $A$ et $B$ sont $-5$, $-9$.
\it Les abscisses de $A$ et $B$ sont $-9$, $+5$.
\end{itemize}
\it Soient sur un axe $(xy)$ les points $A$, $B$, $C$, $D$ d'abscisses $+2$, $+4$, $-1$, $-3$. Vérifier la relation $\overline{AD} = \overline{AB} + \overline{BC}
+\overline{CD}$ ; en est-il ainsi, quelle que soit la disposition des points $A$, $B$, $C$, et $D$ ?
\it Soient quatre points $A$, $B$, $C$, $D$ sur un axe. Démontrer la relation : $$\overline{AB} + \overline{BC} + \overline{CD}
+\overline{DA}=0.$$
\it Calculer le nombre d'années qui s'est écoulé entre les dates suivantes : $51$ ans avant J.-C. et 800 ans après J.-C. Généraliser avec les années $a$ et $b$. Tenir compte du fait qu'il n'y a pas d'année zéro.
\it Soient $A$ et $B$ d'abscisses $+4$ et $-3$. Déterminer les abscisses des points $C$ et $D$ qui partagent $[AB]$ en trois parties égales. Généraliser en désignant par $a$ l'abscisse de $A$ et par $b$ celle de $B$. 
\it Soient $A$, $B$, $C$, et $D$ quatre points d'un axe dont les abscisses sont $+2$, $-3$, $-1$ et $+4$. Vérifier la relation : \[ \overline{DA}.\overline{BC}+\overline{DB}.\overline{CA}+\overline{DC}.\overline{AB}=0.\]
En est-il toujours ainsi lorsqu'on désigne par $a$, $b$, $c$, $d$ les abscisses des points $A$, $B$, $C$ et $D$.
\end{itemize}
 \chapter{Puissances}

\begin{itemize}
\it Calculer les puissances d'exposant $2$, $3$, $4$, $5$ des nombres suivants : 
\[ (+1); \phantom{miou}(+3); \phantom{miou}(+5); \phantom{miou}(-2); \phantom{miou}(-4); \phantom{miou}(-5).\]
\it Calculer : \[ (-11)^3; \phantom{miou}(+9)^2; \phantom{miou}(-10)^5; \phantom{miou}(+7)^4; \phantom{miou}(-13)^4.\]
\it Calculer : \[ \left(-\frac12\right)^5; \phantom{miou}\left(+\frac23\right)^2; \phantom{miou}\left(-\frac34\right)^3; \phantom{miou}(+1,4)^2; \phantom{miou}(-0,25)^2.\]
\it Calculer : \[ 5^3\times 5^2; \phantom{miou}(-7)\times(-7)^2; \phantom{miou}\left(\frac23\right)\times \left(\frac23\right)^3; \phantom{miou}
\left[\left(-\frac14\right)^2\right]^2\left[(-2)^2\right]^3\]
\item Calculer : 
\begin{itemize}
\it $[(-1)(+2)(-3)]^2$.
\it $[(+1)(-3)(-5)]^3$.
\it $[(+2)(-2)(-3)]^2$.
\it $\displaystyle\left[\left(\frac12\right)\left(-\frac13\right)\left(-\frac15\right)\right]^2$.
\it $ \displaystyle\left[(-5)\left(-\frac35\right)
\left(+\frac13\right)\right]^3$.
\it $\displaystyle \left[ (-0,5)(+0,75)\left(+\frac12\right)\right]^3$.
\end{itemize}
\it Calculer : 
\[ (-5)^7\div (-5)^3; \phantom{miaou}
(+0,01)^3\div (+0,01); \phantom{miaou} 
\left(-\frac45\right)^2\div \left(-\frac45\right)^2.\]
\[ 4^3\div 4^6; \phantom{miaou} (-2,5)^2\div (-2,5)^3; 
\phantom{miaou} \left(\frac23\right)^3\div \left(\frac23\right)^6.\]
\it Calculer la valeur de $x^4-x^3+x^2-x+1$. 
\begin{enumerate}
\item pour $x=3$;
\item pour $x=-2$;
\item pour $x=\frac13$;
\item pour $x=-\frac12$.
\end{enumerate}
\it Calculer la valeur de $x^3-3x^2+5x-7$. 
\begin{enumerate}
\item pour $x=4$;
\item pour $x=-1$;
\item pour $x=\frac23$;
\item pour $x=-\frac23$.
\end{enumerate}
\it Calculer la valeur de $\displaystyle\frac{x^3-1}{x^2+1}$. 
\begin{enumerate}
\item pour $x=1$;
\item pour $x=-1$;
\item pour $x=\frac14$;
\item pour $x=-\frac14$.
\end{enumerate}
\it Calculer la valeur de $a^4-4a^3b+6a^2b^2-4ab^3+b^4$. 
\begin{enumerate}
\item pour $a=2$ et $b=-3$;
\item pour $a=-\frac12$ et $b=-\frac12$.
\end{enumerate}
\it Calculer la valeur de $\displaystyle\frac{3xy}{x^2+y^2}-frac{x-y}{x+y}$.
\begin{enumerate}
\item pour $x=1$ et $y=-2$;
\item pour $x=-\frac12$ et $y=+\frac14$.
\end{enumerate}
\it Soient $3$ points $A$, $B$, $C$ d'abscisses $+2$, $-1$, $-3$, sur un axe où $O$ est l'origine des abscisses. Vérifier la relation : 
\[ \overline{OA}^2.\overline{BC}+\overline{OB}^2.\overline{CA}+\overline{OC}^2.\overline{AB} + \overline{AB}.\overline{BC}.\overline{CA} = 0.\]
\it Sur un axe où $O$ est l'origine des abscisses, sont placés trois points $A$, $B$, $M$ d'abscisses $+3$, $-3$, $+1$. Vérifier les relations : 
\[\overline{MA}.\overline{MB}=\overline{MO}^2-\overline{OA}^2.\]
\[ \overline{MA}^2+\overline{MB}^2=2\overline{MO}^2+2\overline{OA}^2.\]
\[\overline{MA}^2-\overline{MB}^2=2\overline{AB}.\overline{OM}.\]
\end{itemize}
 \chapter{Inégalités}

\begin{itemize}
\item Comparer les nombres relatifs suivants : 
\begin{itemize}
\it $+2$ et $-5$; \phantom{miaou}
$+\frac14$ et $-\frac54$; \phantom{miaou}
$+3$ et $0$; ; \phantom{miaou}
$0$ et $-9$. 
\it $+7$ et $+6$; \phantom{miaou}
$-5$ et $-9$; \phantom{miaou}
$-3$ et $-3,5$; ; \phantom{miaou}
$-0,54$ et $-0,57$. 
\it $+\frac34$ et $+\frac23$; \phantom{miaou}
$-\frac35$ et $-\frac57$; \phantom{miaou}
$-3$ et $-\frac{17}3$; ; \phantom{miaou}
$-5$ et $-\frac{15}2$. 
\it $+2,54$ et $+\frac{17}8$; \phantom{miaou}
$-1,3$ et $-\frac97$; \phantom{miaou}
$-\frac{11}{12}$ et $-\frac{12}{13}$; ; \phantom{miaou}
$+\frac{34}{43}$ et $+\frac{35}{44}$. 
\end{itemize}
\it Quels sont les nombres $x$ qui vérifient à la fois les deux inégalités suivantes : $x<3$ et $x<-1$ ? Interprétation graphique : Comparer la position sur un axe, du point $M$ d'abscisse $x$ par rapport aux points d'abscisses $+3$ et $-1$.
\it Quels sont les nombres $x$ qui vérifient à la fois les deux inégalités $x<-1$ et $x > -5$ ? Interprétation graphique. 
\it Existe-t-il des nombres $x$ qui vérifient à la fois les deux inégalités $x>0$ et $x<2$ ? Interprétation graphique.
\it Quels sont les nombres $x$ qui vérifient à la fois :
\begin{enumerate}
\item La double inégalité $2<x<6$ ? 
\item L'une des inégalités $x<3$ ou $x>5$ ? Interprétation graphique.
\end{enumerate}
\it Existe-t-il des nombres $x$ qui vérifient à la fois : \begin{enumerate}
\item La double inégalité $-3<x<-1$ ?
\item L'une ou l'autre des inégalités $x<-4$ ou $x>0$ ? 
Interprétation graphique. 
\end{enumerate}
\it Quels sont les nombres $x$ qui vérifient à la fois les inégalités suivantes : $x>0$; $x<5$, $x<3$ et $x>2$ ? Interprétation graphique.
\it Quels sont les nombres $x$ qui vérifient à la fois : 
\begin{enumerate}
\item La double inégalité $-\frac35<x<+\frac23$ ?
\item L'une ou l'autre des inégalités $x>\frac29$ ou $x<\frac{5}{12}$ ? Interprétation graphique.
\end{enumerate}
\end{itemize}

 \chapter{Expressions algébriques}
 
\begin{itemize}
\item Calculez la valeur numérique des expressions suivantes : 
\begin{itemize}
\it $3a^2-2a+5$ pour $a =  +4$, $-2$, et $-1$.
\it $3x^2-5x+7$ pour $x =  +3$, $+5$, et $-3$.
\it $4x^3-12x^2-4x+7$ pour $a =  +5$, $-3$, et $\frac12$.
\it $\frac{7x^2}2- \frac{8x-6}3 + \frac{3x+7}4$ pour $a =  +3$, $-2$, et $+5$.
\it $\frac{7-3x}{12}+\frac{2(x-2)}3 +\frac54$ pour $a =  +4$, $+2$, et $-3$.
\it $(x^2+1)^2 - x^4 -2x^2$ pour $a =  -1$, $+\frac12$, et $-\frac12$.
\it $(a^2+1)(a^2-1) + 4a^2$ pour $a =  -3$, $+3$, et $+\frac13$.
\it $\frac{(x+y)^2-(x^2+y^2)}{xy}$ pour $x=-5$ et $y=+2$.
\it $\frac{a^2-ab}{a^2-2ab+b^2}$ pour $a=+7$ et $b=-2$.
\end{itemize}
\it Vérifiez que les expressions suivantes donnent la même valeur pour $a=-5$ et $b=3$ : 
\[ a^2 - ab + b^2; \ \ \ \ \ \frac{a^3+b^3}{a+b}; 
\ \ \ \ \ (a+b)^2-3ab\]
\it Vérifiez que les expressions suivantes donnent la même valeur pour $a=+4$ et $b=-1$ : 
\[ (a+b)^2(a-b); \ \ \ \ \ (a^2-b^2)(a+b); 
\ \ \ \ \ \frac{a^4+b^4-2a^2b^2}{a-b}\]
\it Peut-on calculer pour $a=+5$, et $b=-2$ la valeur de l'expression : \[ \frac{a^2+2ab-3b^2}{a(a-1)-5b^2}.\]
\it Peut-on calculer pour $a=+1$, et $b=+2$ la valeur de l'expression : \[ \frac{3a^2+5b^2}{4a^2-b^2}.\]
\item Réduire les monômes suivants et calculer leurs valeurs numériques :
\begin{itemize}
\it $\left(-\frac23\right) a^2x \times (-3y) \times \left(+\frac25\right)$ pour $a=-3$, $x=2$ et $y=-1$. 
\it $xy \times \left(-\frac23\right)x^2\times \frac34a^2$ pour $a=+5$, $x=-2$, et $y=+3$.
\it $\frac27a^2\times\left(-\frac34\right)xy^3\times \left(-\frac25\right)a^2x$ pour $a=+3,5$, $x=+3$ et $y=-2$. 
\it $\left(-\frac35\right)a^2 \times \left(\frac23\right)b^2x \times (-x^4)$ pour $a=4$, $b=-1$, et $x=-2$. 
\it $4x^3 \times (-3y^2) \times \left(-\frac56\right) a^2x^2y^5$ pour $a=-\frac12$, $x=+4$, et $y=\frac32$. 
\end{itemize}
\item Effectuez les sommes de monômes suivantes : 
\begin{itemize}
\it $\frac23ax-\frac12ax+\frac34ax-\frac56ax$.
\it $-\frac35a^2bx+\frac14a^2bx-\frac72a^2bx+\frac1{10}a^2bx$.
\it $-\frac47a^2b^3x+\frac52a^2b^3x-\frac54a^2b^3x$.
\it $\frac34a^2b^3x^4y-\frac23a^2b^3x^4y+\frac14a^2b^3x^4y$. 
\end{itemize}
\it Dire quels sont les degrés des monômes suivants : 
\[ - \frac23a^2x,\phantom{miaou}
\frac34ab^3c^2,\phantom{miaou}
-\frac52a^4b^5y,\phantom{miaou}
-\frac75a^6x^3y^2.\]
\begin{enumerate}
\item Par rapport à la lettre $x$, puis par rapport à $y$. 
\item Par rapport à l'ensemble des lettres $x$ et $y$, puis par rapport à l'ensemble des lettres $a$, $b$, $x$ et $y$.
\end{enumerate}
\end{itemize}
 \chapter{Polynômes}

\begin{itemize}
\item Réduire et ordonner les expressions suivantes: 
\begin{itemize}
\it $-\frac32 x + \frac54 x-3x^2+\frac{x}6-\frac52 x^2+5+4x^2$.
\it $\frac32 x^2+xy+y^2-2xy+\frac{x^2}3-\frac32x^2$.
\it $4a^2-\frac23a-\frac35a^2+\frac13a-5a-\frac2{15}a^2$.
\it $3x^2+\frac45 - \frac53x - 2x^2 -\frac35x^3 + 4 - 2 x^2 + 7x$. 
\it $4x^2 - \frac72 + \frac35x - \frac52x^2 + \frac43 x^3 - 5 + \frac32 x^3 + 7 - 2x$. 
\it $\frac25 a^2b + 3a^3 -4ab^2 + \frac52 a^2b + \frac72b^3 - b^3 + 2ab^2$. 
\end{itemize}
\it Soient les deux polynômes $A= -4x^3 - 2x+2$ et $B= 4x-6x^2 + 5x^3 - 2$. \begin{enumerate}
\item Calculer $A+B$. 
\item Calculer en $x=2$ les valeurs numériques de $A$, $B$ et $A+B$ pour vérifier la réponse précédente. 
\end{enumerate}
\it Mêmes données que l'exercice précédent.
\begin{enumerate}
\item Calculer $A-B$.
\item Calculer en $x=3$ les valeurs numériques de $A$, $B$ et $A-B$ pour vérifier la réponse précédente. 
\end{enumerate}
\item Réduire les expressions suivantes : 
\begin{itemize}
\it$ \left(-5 x^4 + 3 - \frac45 x^3\right) + \left( 
-\frac23 x^3 - 2x\right) - \left( 7x^2 - \frac45 x + 5x^4\right)$.
\it $(12x^3+2x^2-5x+13) + (3x+5-4x^3)- (5x^3-8+2x^2)$.
\it $(a^3-3a^2b+3ab^2-b^3) + (a^3+3a^2b+3ab^2+b^3) - (6ab^2-3a^3)$. 
\end{itemize}
\item Effectuer : \begin{itemize}
\it $(3x-5)+[2x-5 - (3x-2y+4) - (4x-3y-9)]$.
\it $(2x-5y+7) - [(3x+2y-3)-(4x+4y-2)]-[2x-(3y+4)]$. 
\it $[(x-2y+5) - (3x+2y+7)]-[(2x+3)-(4y-2)]$. 
\end{itemize}
\it Soient les polynômes : $A=3x^2-4x+5$, $B=2x^2+5x-4$, $C= 4x^2-x+3$. 
Calculer et réduire : $A+B+C$, $A+B-C$, $A-B+C$, et $-A+B+C$. 
\it Soient les polynômes : $A=5a^2-3ab+7b^2$, $B=6a^2-8ab+9b^2$, $C= 4a^2-3ab-7b^2$. 
Calculer et Réduire : $A-B-C$, $-A-B+C$ et $-A+B-C$.
\it Soient les polynômes : $P=2x^5-3x^2+4x$, $Q=4x^3 - 5x^2 +2x-1$, $R = 4x^5 - 2x^3 + 3x -1$, et $S= 3x^2+2x-5$. 
Calculer et réduire : $(P+Q)-(R+S)$, $(P-Q)+(R-S)$, et $P-Q-R+S$. 
\end{itemize}
 
 \chapter{Multiplication des polynômes}
  
\begin{itemize}
\item Effectuer les produits suivants : 
\begin{itemize}
\it $(3a^2b^3)\left(\frac23ab^5\right)$.
\it $\left(\frac45a^3b^2c\right)\left(-\frac34abc^4\right)$. 
\it $\left(\frac47a^2xy^3\right)\left(-\frac52a^3y^4\right)$.
\it $\left(-\frac34x^2y\right)\left(+\frac35a^3y^5\right)$.
\it $\left(\frac94a^4x^2y^3\right)\left(-\frac43ax^2\right)$. 
\it $\left(\frac{14}3a^2b^3x\right)\left(-\frac67a^2b^5\right)$. 
\it $\left(-\frac72ax^2y\right)\left(-\frac8{15}b^3xy^2\right)\left(\frac5{21}abx^3\right)$.
\it $\left(-\frac23xy^2\right)^2(-4x^2y)$.
\it $\left(\frac5{12}a^4b^2x\right)\left(-\frac27ax^2y^3\right)\left(-\frac{14}5b^2xy^4\right)$.
\it $\left(\frac35x^2y\right)^3\left(-\frac54xy\right)$.
\end{itemize}
\item Calculer : \begin{itemize}
\it $\left(-\frac25ab^3\right)^2$.
\it $\left(\frac53a^2b^3x^4\right)^2$.
\it $\left(-\frac32a^4b^3y^2\right)^3$.
\it $\left(\frac72a^3b^5x^3\right)^2$.
\it $\left(-\frac94a^4b^2x^5\right)^2$.
\it $\left(-\frac65ax^4y^5\right)^3$.
\end{itemize}
\item Effectuer les produits suivants : 
\begin{itemize}
\it $\left(\frac32a^2b-\frac54ab+3a\right)\left(-\frac43a^2b^3\right)$.
\it $\left(\frac54ax^2+\frac32bx-4x\right)\left(-\frac45ax^5\right)$.
\it $\left(\frac25a^2x-3ay-4by\right)(4a^3x^2y)$.
\it $\left(-\frac32x^5+\frac{15}4x^3-\frac25x\right)\left(-\frac{20}3x^4\right)$.
\it $(2x-3y)(4x-2)$.
\it $(2a+3b)(-4a+6b)$.
\it $(-4x+3y+1)(y-3)$. 
\it $(-2a+3b-5)(a-b)$. 
\it $(2x^3-3y-2+5)(x^2-y)$.
\it $(4a^3-5b^4+ab)(a^2-b)$.
\it $(5xy+3x-2y)(2x-y)$.
\it $(-3xy+4x-2y)(x+5)$.
\it $(14a^2b+5a^2-b)(a^2-2b)$.
\it $(7a^3b-4b^2+2a^3)(2a^3+4b^2)$. 
\end{itemize}
\it Soient les expressions littérales $A=-2x^2+3x+5$ et $B=x^2-x+3$. 
\begin{enumerate}
\item Calculer le produit $AB$. 
\item Pour $x=-3$, vérifier le résultat en calculant séparément les valeurs numériques de $A$, $B$ et $AB$. 
\end{enumerate}
\it Soit l'expression littérale $A= x^2-3x+2$. \begin{enumerate}
\item Calculer le carré, puis le cube de $A$. 
\item Vérifier pour $x=-4$, les valeurs de $A$, $A^2$ et $A^3$. 
\end{enumerate}
\item Effectuer les produits suivants.
\begin{itemize}
\it $(2x-7)(-3x+2)$.
\it $(4x^5+7-2x^3)(x^3-2x)$.
\it $(5x^3-2x)(3x-4x^2)$
\it $(2x-7x^2+5x^3)(3x-5x^2+8)$.
\it $\left(-2x+\frac32\right)(4x+3)$.
\it $\left(\frac83x-\frac32x^2+5\right)(4x^3-5x^2+7)$.
\it $(7x^4-2x^3+4x^2)(3x^2-5)$.
\it $(2x^2-4x^3)(x^3-2x)$.
\it $(2x^2-4+2x)(x^2+5-2x)$.
\it $\left(\frac54x^3-2x+\frac12\right)\left(\frac72x^3-\frac23x+x^2\right)$.
\end{itemize}
\item Calculer les produits suivants : 
\begin{itemize}
\it $(2x+3)(3x+2)(x-4)$.
\it $(5x-1)(2x+3)(7+4x)$.
\it $(3x^2-1)(x+1)(x-1)$. 
\it $\left(x-\frac35\right)(5x^2-1)(5x+3)$.
\it $(2x^2+3x-4)^2$. 
\it $(4x^3-7x+2x^2+5)^2$. 
\it $(7x-5)^3$. 
\it $(x^2-x+2)^3$
\end{itemize}
\item Développer et réduire : 
\begin{itemize}
\it $5(3a^2-4b^3)-[9(2a^2-b^3)-2(a^2-5b^3)]$.
\it $3a^2(2b-1)-[2a^2(5b-3)-2b(3a^2+1)]$.
\it $(2a+5b)(3a-2b)-(2a-1)(3a+2b)-(a-2b)(5b-1)$.
\it $(2x-3y)(5x-2y)-(3x-2y)(2x+1)-(5x-y)(3y+1)$.
\it $(ax^2-b)(ax^2-2b)+3b(ax^2-b)+b(b-1)$.
\it $(x-1)(x-2)(x-3)+6(x-1)(x-2)+7(x-1)$.
\it $(x^2+y^2)(x^2-y^2)(x-y)+xy(x^3+y^3)$.
\it $\frac23x^2y\left(2x^2-\frac{y}3\right)-2x^2(2x^2-1)+\left(2x^2-\frac{y}3\right)\left(1-\frac{y}3\right)(2x^2-1)$;
\end{itemize}
\end{itemize}
 \chapter{Identités remarquables}


\begin{itemize}
\item Vérifiez les identités suivantes : 
\begin{itemize}
\it $\frac12(a+b)^2+\frac12(a-b)^2=a^2+b^2$.(identité du parallélogramme).
\it $\left(\frac{a+b}2\right)^2-\left(\frac{a-b}2\right)^2=ab$. (identité de polarisation). 
\it $(a-b)(a^3+a^2b+ab^2+b^3)=a^4-b^4$.
\it $(a+b)(a^3-a^2b+ab^2-b^3)=a^4-b^4$.
\it $(x^2+x+1)(x^2-x+1)= x^4+x^2+1$.
\it $(aa'+bb')^2+(ab'-a'b)^2=(a^2+b^2)(a'^2+b'^2)$.
\it $(x-1)(x+1)(x^2+1)=(x-1)(x^3+x^2+x+1)=x^4-1$.
\it $a(b-c)+b(c-a)+c(a-b)=0$.
\it $a(bz-cy)+b(cx-az)+c(ay-bx)=0$.
\it $(x+y)^3-3xy(x+y)=x^3+y^3$.
\it $(x+y)^3+2(x^3+y^3)=3(x+y)(x^2+y^2)$. 
\end{itemize}
\item Utiliser les identités remarquables du cours pour développer les produits suivants : 
\begin{itemize}
\it $\left(\frac32x^3-\frac25y^2\right)^2$.
\it $\left(\frac43x^5+\frac25y^3\right)^2$.
\it $\left(\frac25x^2-\frac34y\right)\left(\frac25x^2+\frac34y\right)$
\it $\left(\frac23a^2x^3-\frac12by^4\right)\left(\frac23a^2x^3+\frac12by^4\right)$.
\it $(3x+4y-5)(3x+4y+5)$.
\it $\left(\frac23x-\frac45y-1\right)\left(\frac23x+\frac45y+1\right)$.
\it $(3x+4y-2z)^2$. 
\it $\left(\frac52x-\frac34y+z\right)^2$. 
\end{itemize}
\item Développer et réduire : 
\begin{itemize}
\it $(a+b)(a+x)(b+x)-a(b+x)^2-b(a+x)^2$.
\it $bc(b-c)+ca(c-a)+ab(a-b)+(b-c)(c-a)(a-b)$.
\it $(a+b+c)[(a-b)^2+(b-c)^2+(c-a)^2]$.
\it $(b-c)(x-a)^2+(c-a)(x-b)^2+(a-b)(x-c)^2$.
\it $(a+b)^2+(b+c)^2+(c+a)^2-(a+b+c)^2$. 
\it $a^2(a-b)(a-c)+b^2(b-c)(b-a)+c^2(c-a)(c-b)$. 
\end{itemize}
\end{itemize}
 \chapter{Division des monômes et des polynômes}

\begin{itemize}
\item Calculer les quotients de : 
\begin{itemize}
\it $-\frac23a^4x^5y^2$ par $\frac52ax^3y^2$.
\it $\frac35a^4b^2c^5$ par $-\frac34ab^2c^3$.
\it $-\frac{10}7a^5x^2y^7$ par $\frac{4}{21}a^2x^2y^3$. 
\it $-\frac9{20}a^3x^2y^6$ par $-\frac34a^2x^2y$.
\it $3a^5x^4y^3$ par $-\frac94a^4x^2y^3$. 
\it $-\frac72x^3y^5z^2$ par $\frac{21}4x^3y^2z$. 
\end{itemize}
\item Effectuer la division de : \begin{itemize}
\it $14a^5-21a^3+3a^2$ par $\frac72a^2$. 
\it $7a^3x^2-4a^3x+8a^2x^4$ par $-3a^2x$.
\it $-\frac{10}3 a^2bx^3+\frac{15}2ab^3x^4-5abx^2$ par $\frac53abx^2$.
\it $-\frac{21}5a^2b^3x^2y^3+\frac23a^3bx^4y-7a^2bxy^2$
par $-\frac72a^2bxy$.
\end{itemize}
\item Mettre en facteur le monôme de plus haut degré possible dans les polynômes. 
\begin{itemize}
\it $5x^6-7x^3-4x^5+2x^2$.
\it $2a^3x^2-4ax^3+7a^2x^3$.
\it $\frac65a^2b^3x^2-\frac72a^3bx^4+\frac53a^4b^2x^3$.
\it $\frac32a^2b^3x^2y^3+\frac52a^3bx^4y^2-7a^2xy^4$.
\end{itemize}
\item Décomposer en un produit de facteurs les expressions suivantes : 
\begin{itemize}
\it $\frac52x^3y^2-5x^2y^2+5xy^2$.
\it $18abx^2-12abx+2ab$.
\it $\frac45a^3x^2-5a^3y^2$.
\it $\frac34a^2bx^2-\frac{25}3a^2by^2$.
\it $25a^2x^4y^2-4b^2y^2$.
\it $(2x-3)^2-(3x-5)^2$.
\it $(2x+3)^2-4(2x+3)$.
\it $(5x^2+3x-2)^2-(4x^2-3x-2)^2$.
\it $(3x-5)^2-(3x-5)(2x+3)$.
\it $(a^2+b^2-2)^2-(2ab-2)^2$.
\it $a(x^2+1)-x(a^2+1)$.
\it $ab(x^2+y^2)+xy(a^2+b^2)$.
\it $(a^2+b^2-10)^2-(a^2-b^2-8)^2$.
\it $(4a^2+b^2-9c^2)^2-16a^2b^2$.
\end{itemize}
\item Simplifier : 
\begin{itemize}
\it $\frac{7a^2x^5}{b^2x^4}$.
\it $\frac{-2a^3b^2x}{3a^3bx^3}$.
\it $\frac{10a^2x^3y^2}{-4a^4x^3y}$.
\it $\frac{x^2}{x^2-x}$.
\it $\frac{x^2+x^3}{x^3-x}$.
\it $\frac{6x^2-4ax}{9ax-6a^2}$.
\it $\frac{ax+by}{a^2x^2-b^2y^2}$.
\it $\frac{x^3-9x}{3x^2-9x}$.
\it $\frac{x^2+2x+1}{x^2-1}$.
\end{itemize}
\item Calculer les expressions suivantes : 
\begin{itemize}
\it $\frac{x+1}6+\frac{2(2x-1)}{21}-\frac{3x+1}{14}$.
\it $\frac{x+2}5-\frac{4x+1}{15}-\frac{x+1}3$.
\it $\frac1{a(a+1)}-\frac1{a(a-1)}+\frac{2a}{a^2-1}$.
\it $\frac{2}{2a+1}-\frac1{2a-1}+\frac2{4a^2-1}$.
\it $\frac1{x-1}-\frac3{x+1}+\frac{x^2-3}{x^2-1}$. 
\it $\frac{x}{x-1}-\frac2{x+1}-\frac2{x^2-1}$.
\it $\frac2{x+2}-\frac1{x-2}+\frac4{x^2-4}$.
\it $\frac{x-2}{x^2+2x}-\frac1{x+2}+\frac1x$.
\it $\frac{x^3}{x^3-x^2}+\frac{x}{x+1}-\frac{2x}{x^2-1}$.
\it $\frac1x+\frac{x-2}{x^2-4}-\frac2{x^2+2x}$.
\end{itemize}
\item Simplifier les expressions suivantes : 
\begin{itemize}
\it $\frac{x}2\times\frac{x+1}6\times\frac{4x}{x^2-1}$.
\it $\frac{x+3}5\times\frac{x+1}{x^2}\times\frac{x}{(x+3)(x+1)}$.
\it $\frac{a-b}a\times\frac{a^2-ab}5\times\frac{3a}{a^2-b^2}$.
\it $\frac1{a^2-ab}\times\frac4{a^2}\times\frac{a^2-b^2}5$.
\it $\frac{\frac{x}{x-y}-\frac{y}{x+y}}{\frac{x}{x+y}+\frac{y}{x-y}}$.
\it $\frac{\frac{x+1}{x-1}-\frac{x-1}{x+1}}{1-\frac{x-1}{x+1}}$.
\it $\frac{\frac{a}b+\frac{b}a-2}{\frac{a}{b}-\frac{b}{a}}$.
\it $\frac{\frac{a+b}2+\frac{b^2}{a-b}}{\frac{a+b}2-\frac{ab}{a+b}}$.

\end{itemize}
\end{itemize} 
 \chapter{Équation du premier degré à une inconnue}


\begin{itemize}
\item Quel est le degré des équations suivantes ? \begin{itemize}
\it $(4x-1)^2=(2x+3)^2$.
\it $2x(x+5)-(x-3)^2=0$.
\it $(2x-1)^2+(x+3)^2-5(x+7)(x-7)=8$.
\it $(x+1)^3+2(x-1)^3+x^3-3x(x+1)(x-1)=0$.
\end{itemize}
\item Résoudre les équations suivantes : 
\begin{itemize}
\it $5(2x-3) - 4(5x-7)=19-2(x+11)$. 
\it $4(x+3)-7x+17 = 8(5x-1)+166$. 
\it $17-14(x+1)= 13-4(x+1)-5(x-3)$. 
\it $5x+3,5+(3x-4) = 7x-3(x-0,5)$. 
\it $7(4x+3)-4(x-1)=15(x+0,75)+7$. 
\it $17x+15(x-1)=-1-14(3x+1)$. 
\end{itemize}
\item Résoudre les équations suivantes (après développement, les termes en $x^2$ ou $x^3$ se simplifient) : 
\begin{itemize}
\it $(x-1)^2+(x+3)^2= 2(x-2)(x+1)+38$.
\it $5(x^2-2x-1)+2(3x-2)=5(x+1)^2$.
\it $(9x+1)(x-2)=(3x+4)(3x-5)$.
\it $7(3-2x)-5x(2x-1)=(5x+3)(3-2x)$. 
\it $(3x-1)^2-(2x+3)^2+7 = (2x+1)(2x-1) + x(x+7)$.
\it $(x+2)^3+(x-2)^3+(x+1)^3=3(x+1)(x+2)(x-2)$.
\end{itemize}
\item Résoudre les équations suivantes : 
\begin{itemize}
\it $\frac52x+3 - \frac{7x}4 = x + \frac94$. 
\it $\frac{3x}7 - \frac{2x}{15} + 3 = \frac{x}3+\frac{13}3$.
\it $x+\frac12-\frac{x}6 = 16 - \frac{2x}9 + \frac13$. 
\it $\frac{7x}4-2-\frac{x}2=\frac{2x}{13}-\frac{85}{52}$. 
\it $\frac{2x}3+4 -\frac{2x}5 = \frac{x}2-\frac{x}3+3,5$. 
\it $\frac{x}6-1=\frac{x}4-\frac{x}3-1$. 
\end{itemize}
\item Résoudre les équations suivantes (on se ramènera au cas entier en multipliant par un nombre opportun).
\begin{itemize}
\it $\frac{x+5}4 - \frac{x-3}6 = \frac{x}3$. 
\it $\frac{3x-7}2 + \frac{x+1}3 = -16$. 
\it $x-\frac{x+1}3 = \frac{2x+1}5$. 
\it $\frac{7-3x}{12} + \frac34 = 2(x-2) + \frac{5(5-2x)}6$. 
\it $\frac{x}5 - \frac{3x-1}6 + \frac{3-x}4 = 0$.
\it $\frac{3(x+3)}4 + \frac12 = \frac{5x+9}3 - \frac{7x-9}4$. 
\it $\frac{2x-7}5 + \frac{x+11}2 = -4$. 
\it $\frac{2x-3}3 - \frac{x-3}6 = \frac{4x+3}4 - 17$. 
\it $\frac{5x-3}4 - \frac{7x-5}9 = \frac{x+19}6$. 
\it $\frac{5x+1}8-\frac{x-1}3 = \frac{4(2x-3)}9$. 
\it $\frac{2x-1}3 - \frac{5x+2}7 = x +13$. 
\it $\frac{8x+2}5 - \frac{x-11}7 = \frac{5x-3}2 - \frac{3x-1}4$.
\it $\frac{2x-7}9 - \frac{x-5}6 = \frac{x-9}8$. 
\it $\frac{5x+7}4 - \frac{3x+5}8 = \frac{4x+9}5 - \frac{x-9}3$. 
\it $\frac{5x+6}7 - \frac{3x+1}4 = \frac{x+16}5$. 
\it $\frac{4x+7}5 -\frac{x-5}6 = \frac{2x+14}3 - \frac{2x-7}9$. 
\end{itemize}
\item Résoudre les équations suivantes (après multiplication et développement, les termes en $x^2$ disparaissent) : 
\begin{itemize}
\it $\frac{(x-1)(x+5)}3 - \frac{(x+2)(x+5)}{12} = \frac{(x-1)(x+2)}4$.
\it $\frac{(x+1)^2}3 + \frac{(x-2)(x-3)}2 = \frac{(5x-1)(x-4)}6 + \frac{28}3$. 
\it $\frac{(3x+1)(3x-1)}9 - \frac{(x-5)(x+1)}2 = \frac{(9x-1)(x+3)}{18} + \frac89$. 
\it $\frac{(4x+7)^2}4 - \frac{(5x-1)^2}7 = \frac{(8x-3)(3x+4)-79x}{56}$. 
\it $\left(x-\frac83\right)(x+0,75) = (x+4,5)(x+1,5) - \frac{145}3$. 
\it $\frac{(x-5)^2}5 + \frac{(x+3)^2}3 = \frac{(3x+1)(3x-1)-x(x+1)}{15}$. 
\it $\left(3x-\frac45\right)\left(5x+\frac23\right)
=15 (x-1)(x+1)+\frac7{15}$. 
\end{itemize}
\end{itemize}

 
 \chapter{Équations se ramenant au premier degré}
\begin{itemize}
\item Résoudre les équations produits :
\begin{itemize}
\it $(x-1)(x+2)(x-3)=0$. 
\it $(x-3)(x-4)(x-5)=0$. 
\it $(2x+1)(x+1)(4x-3) = 0$.
\it $(2x+1)(x+4)(3x+1)=0$. 
\it $x(5x+1)(4x-3)(3x-4)=0$. 
\it $5x(3x-7)=0$. 
\end{itemize}
\item Résoudre en factorisant pour faire apparaître une équation produit : 
\begin{itemize}
\it $x^2-3x=0$. 
\it $5x^2+8x=0$. 
\it $4x^2-\frac{7x}3=0$. 
\it $\frac{x^2}5+x=0$.
\it $-\frac{3x^2}5 + x =0$. 
\it $-\frac{5x^2}7 - \frac{3x}4=0$. 
\it $x(x+1)=x+1$.
\it $(4x-1)(x-3)=(x-3)(5x+2)$.
\it $(x+3)(x-5)+(x+3)(3x-4)=0$.
\it $5(x+1)(x+2)(x-3)=4(x+1)(x+2)(x-4)$. 
\end{itemize}
\item Résoudre les équations suivantes, en utilisant des identités ou des factorisations pour se ramener à une équation produit : 
\begin{itemize}
\it $(x+5)(4x-1)+x^2-25=0$.
\it $(x+4)(5x+9)-x^2+16=0$.
\it $x^2-9=0$.
\it $5x^2-125=0$.
\it $4x^2-49=0$. 
\it $x^2-100=0$.
\it $x^2=81$.
\it $9x^2=64$. 
\it $(x+1)^2-(2x-5)^2=0$.
\it $(2x+7)^2-(4x-9)^2=0$.
\it $(5x+1)^2=(x-1)^2$.
\it $(3x+1)^2=(x-4)^2$.
\it $4(x+1)^2-9(x-1)^2=0$.
\it $(x+7)^2-81(x-5)^2=0$.
\it $5x^3-5x=0$.
\it $(x+1)(x-1)^2-(x+1)(x-2)^2=0$.
\it $3x^2-12x=0$. 
\it $(3x+1)(x-3)^2=(3x+1)(2x-5)^2$. 
\it $7x^3-175x=0$. 
\it $(x+5)(3x+2)^2=x^2(x+5)$. 
\end{itemize}
\item [Terminer de recopier 662s]
\end{itemize}
 \input{20Problemes}
 \chapter*{Exercices de révision}

\section*{Calcul algébrique}


\begin{itemize}
\it On donne sur un axe quatre points $A$, $B$, $C$, $D$. Démontrer la relation : $\overline{AC} + \overline{BD} = \overline{AD} + \overline{BC}$.
\it Soient sur un axe deux points $A$ et $B$ d'abscisses $a$ et $b$. Calculer les abscisses des points $M$ et $N$ qui partagent $[AB]$ en $3$ parties égales.
\it On donne sur un axe trois points $A$, $B$ et $C$ d'abscisses $a$, $b$ et $c$. Trouver l'abscisse du point $G$ tel que $\overline{GA}+\overline{GB}+\overline{GC}=0$.
\it Soient sur un axe $4$ points, $A$, $B$, $C$ et $D$. Démontrer la relation : \[ \overline{AB}.\overline{CD} + \overline{BC}.\overline{AD}+\overline{CA}.\overline{BD}=0.\]
\it On donne sur un axe $3$ points, $A$, $B$ et $M$. On désigne par $O$ le milieu de $[AB]$. Démontrer les relations : 
\[ \overline{MA} + \overline{MB} = 2 \overline{MO}.\phantom{miaoumiaoumiaou} \overline{MA}.\overline{MB}=
\overline{MO}^2-\frac{\overline{AB}^2}4.\]
\it Soient deux points $A$ et $B$ d'un axe, $O$ le milieu de $[AB]$ et $M$ un point quelconque de l'axe.
Démontrer les relations : 
\[ \overline{MA}^2 + \overline{MB}^2 = 2 \overline{MO}^2+\frac{\overline{AB}^2}2.\phantom{miaoumiaoumiaou} \overline{MA}^2-\overline{MB}^2=2\overline{AB}.\overline{OM}\]
\item Effectuer les produits suivants : \begin{itemize}
\it $\left(\frac34 a^2x^3y^2\right)\left(-\frac29 bx^2y^4\right)(-10a^2b^3x^4y^2)$.
\it $(2x+3)(x-5)(3x+1)$.
\it $(2x^3-3x^2+4x)(4x^2+6x+8)$.
\it $(2x-3y)^2(x-y+1)$.
\end{itemize}
\item Calculer les expressions suivantes : \begin{itemize}
\it $x(x+1)(x+2)-3(x-3)(x+2)+2(x-6)$.
\it $(2x+3y)(2y-1)-(2x-y)(5x-1)+(2x-3y)(5x+2y)$.
\it $(5x+3y)(3x-2y)+(5x+2y)(3y+1)-(5x-2y)(3x+1)$.
\it $(ax-y)(ax-2y)+3y(ax-y+1)-y(8y+3)$.
\end{itemize}
\item Établir les identités suivantes : 
\begin{itemize}
\it $(a+b-c)^2+(a-b+c)^2=2a^2+2(b-c)^2$.
\it $(a+b+c)^2-(a-b-c)^2=4a(b+c)$.
\it $(a+b)^3+(a-b)^3 = 2a(a^2+3b^2)$.
\it $(ab+1)^2+(a-b)^2=(a^2+1)(b^2+1)$.
\it $(ab+1)^2-(a+b)^2=(a^2-1)(b^2-1)$.
\end{itemize}
\item Décomposer en un produit de facteurs les expressions suivantes : 
\begin{itemize}
\it $12a^2b-3b^3$.
\it $48x^2y^2-27y^4$.
\it $x^2+6x+9$. 
\it $3x^2-30x+75$.
\it $(2x+3)^2-49$.
\it $x^2-10x+21$.
\it $(2a-b)^2-a^2$.
\it $(2a+b)^2-(a+2b)^2$.
\it $(ax+1)^2-(x+a)^2$. 
\it $(ax+b)^2+(bx-a)^2$.
\end{itemize}
\item Simplifier les expressions suivantes : 
\begin{itemize}
\it $\frac{(a+b)^2-c^2}{(a+c)^2-b^2}$.
\it $\frac{(3a+2b)^2-(a+2b)^2}{a^2-b^2}$.
\it $\frac{x^2-2x+1}{x^2-1}$.
\it $\frac{4x^2-9}{4x^2+8x-21}$.
\it $\frac{a+1}{a-1}+\frac{a-1}{a+1}-\frac{4a}{a^2-1}$.
\it $\frac2{2a+b}-\frac1{2a-b}+\frac{2b}{4a^2-b^2}$.
\it $\frac{x^2}{x-1}+\frac1{x+1}-\frac{2x}{x^2-1}$.
\it $\frac{x+1}x-\frac2{x^2+2x}+\frac{x-1}{x+2}$.
\it $\frac{\frac{x+1}x-\frac{x}{x+1}}{1+\frac{x}{1+x}}$.
\it $\frac{\frac{a-b}b+\frac{2b}{a+b}}{\frac{a+b}a-\frac{2b}{a+b}}$.
\end{itemize}
\it Soit les rapports égaux $\frac{a}{a'}=\frac{b}{b'}=\frac{c}{c'}$. Démontrer que chacun de ces rapports est égal à $\frac{5a-3b+2c}{5a'-3b'+2c'}$.
\it Soit l'égalité $\frac{a}b=\frac{c}d$. Démontrer que $\frac{a^2+b^2}{c^2+d^2} = \frac{ab}{cd}$.
\it Démontrer que si $a$, $b$, $c$, $d$ vérifient la relation $(ad+bc)^2=4abcd$, ces quatre nombres vérifient l'égalité $ad=bc$. 
\it Démontrer que si $(a^2+b^2)(c^2+d^2) = (ac+bd)^2$, les quatre nombres $a$, $b$, $c$ et $d$ vérifient l'égalité $ad=bc$.
\end{itemize}

\section*{Équations du premier degré}
\begin{itemize}
\item Résoudre les équations suivantes : 
\begin{itemize}
\it $\frac{x-1}3+\frac{4(2x-3)}9 = \frac{5x+1}8$.
\it $\frac{x+5}6 + \frac{x+9}8 = \frac{2x+7}9$.
\it $\frac{5x-9}3 + \frac{7x-5}9 = \frac{3x-143}4$.
\it $\frac{13x-16}{15}+\frac{x-32}{35} = \frac{295-x}21$.
\it $\frac{3x-2}{18}+\frac{2x-9}{45}=\frac{x+3}{10}$.
\it $\frac{4x+9}{77}-\frac{5x+12}{14}=\frac{11x+28}{22}$.
\it $\frac{x+5}4 -\frac{x-3}6 = \frac{x}3$.
\it $-\frac{x+1}3 = \frac{3x+4}5$.
\it $\frac{5x-7}4 + \frac{3x-5}8 = \frac{9x-4}5$.
\it $\frac{4x-1}3+\frac{2x-1}2=\frac{10x+63}7$.
\it $\frac{x-8}5-\frac{x-23}{15} = \frac{43-x}{18} $.
\it $\frac{13x-34}{21} - \frac{3x-7}{56} = \frac{19-7x}{24}$. 
\end{itemize}
\item Résoudre les équations suivantes : 
\begin{itemize}
\it $\frac{3x-1}7 + \frac{3x-4}{5x-3} = \frac{6x+5}{14}$. 
\it $\frac{2x-7}{2(x-1)} = 1+ \frac1{4-x}$.
\it $\frac3{x-3} + \frac4{x-5} = \frac7{x-4}$.
\it $\frac{2x+3}{5x-3} = \frac{3}{2(2x-3)} + \frac25$.
\it $\frac{9x+2}{6(x+1)} - \frac{5x+8}{2x+1} = - \frac{x+5}{x+1}$. 
\it $\frac5{2(3x+2)} - \frac1{2x+3} = \frac1{3x-8}$. 
\end{itemize}
\item Résoudre les équations suivantes : 
\begin{itemize}
\it $(2x-3)(4x-12)=0$.
\it $x(x-7)(3x-2) = 0$.
\it $x^2-(4x+5)^2 = 0$. 
\it $(2x-3)^2- (x+6)^2=0$.
\it $\frac4{3-2x}+\frac9{x+3} = \frac{13}{2x+3}$.
\it $\frac4{x+10}+\frac5{x-20} = \frac9{x+20}$.  
\end{itemize}
\item Résoudre les équations suivantes où le nombre $m$ est supposé connu : 
\begin{itemize}
\it $\frac{8x-5m}7 - \frac{5x-11m}5 = \frac{78 - x}{35}$. 
\it $m(x-m)=x-1$.
\it $\frac{4x+m}5 - \frac34(x-1) = \frac{2x-m+3}{10}$.
\it $x+m^2+1=m(x+2)$.
\end{itemize}
\it Soit l'équation : \[ \frac{5x-m}3 - \frac{4x+m}5
= \frac{7(3x+4)}{14}\]
dans laquelle $m$ est un nombre connu.
\begin{enumerate}
\item Résoudre cette équation.
\item Déterminer $m$ de façon que $x=1$.
\item Calculer les valurs de $m$ pour lesquelles $x>3$.
\end{enumerate}
\end{itemize}
\section*{Problèmes du premier degré}
\begin{itemize}
\it La somme de deux nombres est $324$. En ajoutant $26$ à chacun d'eux, l'un devient le triple de l'autre. Quels sont ces deux nombres ? 
\it La différence de deux nombres est $1~032$. Le quotient entier de ces deux nombres est $13$ et le reste de leur division est $48$. Quels sont ces deux nombres ? 
\it Deux cyclistes partent en même temps de deux villes A et B distantes de $200$ km et vont à la rencontre l'un de l'autre. Ils se rencontrent au bout de $4$ heures. Si le cycliste qui part de A était parti une demi-heure avant l'autre, la rencontre aurait eu lieu $3$ heure $48$ minutes après le départ du deuxième cycliste. Quelle est la vitesse de chacun d'eux ? 
\it Un train a mis $36$ secondes à passer devant un observateur immobile. Sa longueur est $300$ m. Quelle est sa vitesse ? Un second train met $24$ secondes pour croiser le premier et passe devant l'observateur immobile en $18$ secondes. Trouver la longueur et la vitesse du second train. 
\it Deux cyclistes sont séparés par une distance de $54$ km. S'ils allaient à la rencontre l'un de l'autre, ils se rencontreraient au bout d'une heure. S'ils allaient dans le même sens, ils se rejoindraient au bout de $9$ heures. Quelle est la vitesse de chaque cycliste ? 
\it Deux capitaux ont pour somme $3~000$ F.L'un est placé à 6\% et l'autre à 5\%. La somme de leurs intérêts annuels est 162 F. Quelle est la valeur de chaque capital ? 
\it Deux lingots d'or sont l'un au titre de 0,95, l'autre au titre de 0,8. On les fond ensemble en ajoutant 3 kg d'or pur. Le lingot ainsi obtenu est au titre de 0,906 et pèse 37,5 kg. Quel est le poids respectif des lingots primitifs. 
\it On veut obtenir 450 g d'un alliage d'or au titre de 0,83 en fondant des lingots aux titres respectifs de 0,8 et 0,9. Quels doivent être les poids de ces deux lingots ? 
\it Un épicier achète 100 kg de café, une partie à 9,60 F le kilogramme, l'autre à 11,20 F le kilogramme. Il revend le tout pour 1~209,60 F avec un bénéfice de 20\% sur le prix d'achat. Quels sont les poids respectifs de chaque qualité de café ? 
\it Un terrain rectangulaire a 414 m de périmètre. Si la longueur augmentait du tiers de sa valeur et si la largeur diminuait du tiers de la sienne, le périmètre augmenterait de 26 m. Quelles sont les dimensions de ce  terrain ? 
\it La longueur d'un rectangle est supérieure à sa largeur de 21 m. Si on augmentait ses dimensions de 4 m, la surface augmenterait de 492 $m^2$.Trouver les dimensions de ce rectangle. 
\it Deux capitaux sont tels que le second surpasse de 2~300 F les $\frac34$ du premier. En les plaçant à 3,6\%, le premier pendant 7 mois, le second pendant 5 mois, l'intérêt du premier dépasse de 207,30 F l'intérêt du second. Calculer les deux capitaux. 
\it Deux villes A et B sont distantes de 600 km. La tonne de charbon coûte 120 F en A et 144 F en B; et le transport coûte 0,15 F par tonne et par kilomètre. Trouver le point situé entre A et B où le charbon revient au même prix, soit qu'il vienne de A, soit qu'il vienne de B. 
\it Une somme de 2~170 F est partagée entre deux personnes. La première ayant dépensé les $\frac57$ de sa part et la deuxième les $\frac25$ de la sienne, il leur reste la même sommme. Trouver ces deux parts. 
\it Deux cyclistes partent en même temps de deux villes A et B et vont à la rencontre l'un de l'autre. Celui qui part de A fait 6 km à l'heure de plus que l'autre et s'arrête un quart d'heure après chaque heure de marche. L'autre, qui part de B, n'a qu'un seul arrêt de 12 minutes. Les cyclistes se rencontrent au bout de 5 heures au milieu de [AB]. Trouver la distance AB et la vitesse de chaque cycliste. 
\it Partager une somme de 41~340 F proportionnellement aux trois nombres 11, 13 et 15. 
\it Les fortunes de 3 personnes sont proportionnelles aux nombres 2, 3 et 4. En additionnant la première, le double de la seconde et le triple de la troisième, on trouve 1~000 F. Quels sont ces trois fortunes ? 
\it Un cultivateur a vendu 2~880 F sa récolte de blé et 1~440 F sa récolte d'avoine. Le poids du blé surpasse de 16 quintaux celui de l'avoine. Sachant que le prix du quintal de blé est les $\frac85$ de celui du quintal d'avoine : 
\begin{enumerate}
\item Trouver les quantités vendues de chacune des deux cérales. 
\item Quel est le prix du quintal de chacune d'elles ?
\end{enumerate}
\it Un automobiliste part à 8h40 d'une ville A et se rend à une ville B située à 39 km de A. Au retour, il fait un détour de 21 km, mais sa vitesse horaire surpasse de 10 km à l'heure sa vitesse à l'aller, si bien que la durée du trajet retour est égale aux $\frac43$ de la durée du trajet aller.\begin{enumerate}
\item Calculer la vitesse de l'automobiliste à l'aller et au retour.
\item Il s'est arrêté 1h 26 min en B. Quelle est l'heure de son retour en A ? 
\end{enumerate}
\it Un chemisier a acheté deux lots de chemises identiques. Il vend le premier lot pour 576 F en réalisant un bénéfice de 2,40 F l'unité. Il vend une partie du second lot pour 180 F avec un bénéfice de 3F et le reste pour 252 F avec un bénéfice de 2 F par chemise. Sahcant que le nombre de chemises du deuxième lot est les $\frac34$ de celui du premier lot : \begin{enumerate}
\item Trouver le prix d'achat d'une chemise.
\item Trouver le nombre de chemises vendues à chaque fois.
\end{enumerate}
\it Un cycliste effectue un trajet comprenant 36 km de terrain plat, 24 km de montées et 48 km de descentes. Dans les montées, sa vitesse horaire moyenne diminue de 12 km et en descente, elle augmente de 15 km. Sachant que al durée du trajet en terrain plat est le tiers de la durée totale du trajet : \begin{enumerate}
\item Trouver la vitesse horaire du cycliste en terrain plat.
\item Quelle est la durée totale du trajet et la vitesse moyenne réalisée ?
\end{enumerate}
\it Un automobiliste part à 8h20 d'une localité A pour une ville B située à 192 km de A et où il doit arriver à une heure déterminée. Il calcule la vitesse horaire moyenne qu'il doit réaliser, mais arrivé à mi-route, il s'aperçoit que sa vitesse horaire a été inférieur de 12 km à la vitesse prévue. Il accélère son allure et arrive à l'heure fixée, en effectuant dans la deuxième partie de son parcours une vitesse horaire supérieure de 18 km à la vitesse prévue. \begin{enumerate}
\item Calculer la vitesse moyenne réalisée sur le parcours entier, puis l'heure d'arrivée. 
\item Quel était le retard à mi-route de l'automobiliste ?
\end{enumerate}
\end{itemize}


\part{Géométrie}
\setcounter{n}{1}
 \chapter{Rappels de définitions}

\begin{itemize}
\it On considère sur une droite $(xy)$ trois points $O$, $A$, $B$ dans cet ordre et on construit le milieu $M$ de $[AB]$. Soit $OA=a$ et $OB=b$. \begin{enumerate}
\item Calculer à partir de $a$ et $b$ les longueurs 
$AM$, $MB$ et $OM$. 
\item Comment faut-il modifier les résultats si $O$ est entre $A$ et $M$ ? 
\end{enumerate}
\it Soit $B$ un point du segment $[AC]$. On désigne par $M$ le milieu de $[AB]$, par $N$ celui de $[BC]$ et on pose $AB = a$ et $BC = b$. \begin{enumerate}
\item Calculer à partir de $a$ et de $b$ la longueur $MN$. 
\item Soit $P$ le milieu de $[AC]$. Calculer $MP$ et $PN$. 
\end{enumerate}
\it Deux angles $\widehat{AOB}$ et $\widehat{BOC}$ sont adjacents et $[OM)$ est la bissectrice de l'angle 
$\widehat{BOC}$. \begin{enumerate}
\item Construire la figure sachant que $\widehat{AOB}=48^o$ et $\widehat{AOC}= 112^o$. Calculer la mesure de
l'angle $\widehat{AOM}$. 
\item Si $\alpha$ et $\beta$ désignent les mesures des angles $\widehat{AOB}$ et $\widehat{AOC}$, montrer que l'on a toujours $\widehat{AOM}= \frac12(\alpha+\beta)$.  
\end{enumerate}
\it Soient $[OM)$ et $[ON)$ les bissectrices de deux angles adjacents $\widehat{AOB}$ et $\widehat{AOC}$. 
\begin{enumerate}
\item Construire la figure pour $\widehat{AOB}= 64^o$ et $\widehat{AOC}= 108^o$.
\item Calculer la mesure de l'angle $\widehat{MON}$. Comparer le résultat à la mesure de $\widehat{BOC}$. 
\item Énoncer un théorème donnant la mesure de l'angle des bissectrices de deux angles adjacents de mesures $\alpha$ et $\beta$.
\end{enumerate}
\it On construit un angle $\widehat{AOB}= 72^o$, puis les angles droits $\widehat{AOA'}$ et $\widehat{BOB'}$ adjacents au premier. 
\begin{enumerate}
\item Calculer la mesure de l'angle $\widehat{A'OB'}$. Comment sont les deux angles $\widehat{AOB}$ et $\widehat{A'OB'}$ ? 
\item Montrer que les bissectrices $[OM)$ et $[ON)$ des angles $\widehat{AOB}$ et $\widehat{A'OB'}$ sont alignées.
\item Calculer l'angle des bissectrices des angles droits $\widehat{AOA'}$ et $\widehat{BOB'}$.
\end{enumerate}
\it Soit un angle $\widehat{AOB}$ de $54^o$. On construit les angles droits $\widehat{AOA'}$ et $\widehat{BOB'}$ non adjacents à l'angle $\widehat{AOB}$.\begin{enumerate}
\item Montrer que les angles $\widehat{AOB}$ et $\widehat{A'OB'}$ sont supplémentaires et qu'ils ont même bissectrice $[OM)$.
\item Démontrer que les angles $\widehat{AOB'}$ et $\widehat{A'OB}$ sont égaux. Calculer l'angle de leurs bissectrices.
\end{enumerate}
\it On construit trois angles successivement adjacents 
$\widehat{AOC} = 48^o$, $\widehat{COD}= 72^o$ et $\widehat{DOB}= 34^o$. Puis on trace les demi-droites $[OM)$, $[ON)$, $[OP)$ et $[OQ)$ bissectrices des angles $\widehat{AOC}$, $\widehat{AOD}$, $\widehat{BOC}$ et $\widehat{BOD}$. \begin{enumerate}
\item Calculer les angles $\widehat{MON}$ et $\widehat{POQ}$ et comparer leur valeur à celle de l'angle $\widehat{COD}$.
\item Montrer que les angles $\widehat{MOQ}$ et $\widehat{NOP}$ ont même bissectrice. 
\end{enumerate}
\it On considère sur un cercle les arcs consécutifs $\arc{AB} = 90^o$, $\arc{BC} = 68^o$ et $\arc{CD} = 90^o$. \begin{enumerate}
\item Calculer la mesure de l'arc $\arc{DA}$ qui termine le cercle. 
\item Construire au rapporteur le milieu $M$ de $\arc{BC}$ et le milieu $N$ de $\arc{DA}$. 
\item On plie la figure suivant la droite $(MN)$. Qu'en déduisez-vous pour le diamètre $[MN]$ et les cordes $[BC]$ et $[DA]$ ? 
\end{enumerate}
\it Soient $I$ et $J$ deux points d'une droite $(xy)$ et $A$ et $B$ deux points extérieurs situés d'un même côté de cette droite. Les cercles de centre $I$ et $J$ passant par $A$ se recoupent en $A'$, les cercles de centres $I$ et $J$ passant par $B$ se recoupent en $B'$. 
\begin{enumerate}
\item Montrer que si on plie la figure suivant $(xy)$, le point $A$ vient en $A'$ et le point $B$ en $B'$.
\item Montrer que $(xy)$ est médiatrice de $[AA']$ ainsi que de $[BB']$. 
\item La droite $(AB')$ coupe $(xy)$ en $M$. Montrer que $A'$, $M$ et $B$ sont alignés.
\end{enumerate}
\it \begin{enumerate}
\item Construire d'un même côté d'une droite $(xy)$ un polygone $F$ de six à huit côtés présentant un angle rentrant, puis le polygone $F'$ symétrique de $F$ par rapport à $(xy)$. 
\item Vérifier que les prolongements de deux côtés homologues $[BC]$ et $[B'C']$ par exemple se coupent sur $(xy)$ en un même point $I$. Comparer les angles $\widehat{BIx}$ et $\widehat{B'Ix}$.
\end{enumerate}
\it Soit un contour polygonal $ABCDEA'$ situé d'un même côté de la droite $(AA')$ et présentant un angle rentrant en $C$ ou en $D$. On désigne par $O$ le milieu de $[AA']$. \begin{enumerate}
\item Construire le contour $A'B'C'D'E'A$ symétrique du précédent par rapport au point $O$. Comparer les angles 
$\widehat{BC'D}$ et $\widehat{B'CD'}$. ainsi que les longueurs $BC'$ et $B'C$. 
\item Les segments $[BC']$ et $[CD']$ se coupent en $M$. Soit $M'$ le symétrique de $M$ par rapport à $O$. Démontrer que $M'$ se trouve à la fois sur $(CB')$ et sur $(DC')$. 
\end{enumerate}
\it On considère un angle droit $\widehat{AOB}$ et on dessine une ligne $F$, brisée ou comprenant des arcs de cercle, joignant $A$ à $B$, notée $AMNPQB$.\begin{enumerate}
\item Construire la figure $F_1$ et la figure $F_3$,
symétriques de $F$ par rapport à $(OA)$ et à $(OB)$.
Puis la figure $F_2$ symétrique de $F$ par rapport à $O$.
\item Montrer que $F_2$ et $F_1$ sont symétriques par rapport à la droite $(OB)$, tandis que $F_2$ et $F_3$ sont symétriques par rapport à $(OA)$, et que $F_1$ et $F_3$ sont symétriques par rapport à $O$. 
\item Quels sont les éléments de symétrie (centre et axes) de la figure totale formée par $F$, $F_1$, $F_2$ et $F_3$ ?
\end{enumerate}
\end{itemize}
 \chapter{Triangles isocèles}
\begin{itemize}
\it Soit un triangle $ABC$ isocèle en $A$. On prolonge $[BA)$ d'une longueur $AD=BA$. \begin{enumerate}
\item Montrer que le triangle $ACD$ est isocèle. 
\item Montrer que l'un des angles du triangle $BCD$ est égal à la somme des deux autres. 
\item Si dans un triangle $ABC$, l'angle en $A$ est égal à la somme des angles en $B$ et $C$, montrer que la médiane $AM$ est égale à la moitié du côté $[BC]$.
\end{enumerate}
\it On construit un cercle de centre $O$ et deux diamètres $[AC]$ et $[BD]$.\begin{enumerate}
\item Démontrer que le quadrilatère $ABCD$ a ses quatres angles égaux. 
\item Quels sont les axes de symétrie de ce quadrilatère ? En déduire l'égalité de $AB$ et $CD$, ainsi que de $AD$ et $BC$.
\end{enumerate}
\it Dans un triangle quelconque $ABC$, la bissectrice extérieure de l'angle en $A$ coupe en $D$ le prolongement $[Cx)$ du côté $[BC]$. On prolonge $[BA]$ d'une longueur $AE=AC$. \begin{enumerate}
\item Que représente la droite $(DA)$ pour le segment $[CE]$, pour le triangle $BDE$ ? 
\item Calculer la mesure de l'angle $BED$ sachant que $\widehat{ACB}=84^o$.
\end{enumerate}
\it On prolonge d'une même longueur et dans le même sens de parcours, les trois côtés d'un triangle équilatéral $ABC$. On obtient les points $A'$ sur le prolongement de $[BC]$, $B'$ sur celui de $[CA]$, et $C'$ sur celui de $[AB]$.\begin{enumerate}
\item Montrer que l'on peut par glissement amener le 
calque de la figure obtenue $ABCA'B'C'$ sur $BCAB'C'A'$ ou $CABC'A'B'$.
\item En déduire que sur les trois triangles $A'B'C$, $B'C'A$ et $C'A'B$ sont directement égaux. Quelle est la nature du triangle $A'B'C'$ ?
\end{enumerate}
\it Soit un quadrilatère convexe $ABCD$ dans lequel $AB=BC$ et $\widehat{BAD}=\widehat{BCD}$.
\begin{enumerate}
\item Montrer que les angles $\widehat{DAC}$ et $\widehat{DCA}$ sont égaux. Comparer $DA$ et $DC$.
\item Que représente $(BD)$ pour le segment $[AC]$ et pour les angles $\widehat{ABC}$ et $\widehat{ADC}$ ?
\end{enumerate}
\it Démontrer que si dans un triangle $ABC$ l'une des conditions suivantes est remplie, le triangle est isocèle. \begin{enumerate}
\item La hauteur $(AH)$ est en même temps bissectrice de l'angle $\widehat{A}$.
\item La hauteur $(AH)$ est en même temps médiane relative à $[BC]$. 
\item La médiatrice de $[BC]$ passe par le sommet $A$.
\end{enumerate}
\it Soit un triangle $ABC$ dans lequel la médiane $(AM)$ est bissectrice de l'angle intérieur $\widehat{BAC}$. On prolonge $[AM]$ d'une longueur $MD=AM$.\begin{enumerate}
\item Montrer que les triangles $MAB$ et $MDC$ sont symétriques par rapport à $M$. En déduire que les angles $\widehat{MAB}$ et $\widehat{MDC}$ sont égaux ainsi que les segments $[AB]$ et $[DC]$. 
\item Nature du triangle $CAD$ ? Comparer $AB$ et $AC$.
\item Énoncer le théorème qui en résulte.
\end{enumerate}
\it Dans un triangle isocèle $ABC$ de base $[BC]$, on mène la médiane $(BM)$. Soit $N$ le milieu de $[AB]$ et $G$ le point de rencontre de $[BM]$ avec la hauteur $(AH)$.\begin{enumerate}
\item On retourne le triangle $ABC$ sur lui-même en amenant $B'$ calque de $B$ en $C$ et $C'$ calque de $C$ en $B$. Où viennent les calques $M'$ et $N'$ de $M$ et de $N$ ? 
\item Montrer que les points $C$, $G$ et $N$ sont alignés et que les triangles $GBC$ et $GMN$ sont isocèles. Comparer $BM$ et $CN$. 
\item Que représente $(AH)$ pour le segment $[MN]$ ?
\end{enumerate}
\it On considère un triangle $ABC$ isocèle en $A$. La bissectrice intérieure de l'angle en $B$ coupe le côté $[AC]$ en $D$ et la hauteur $(AH)$ en $I$. \begin{enumerate}
\item Nature du triangle $IBC$ ? Montrer que $[CI)$ est la bissectrice intérieure de l'angle en $C$. 
\item La bissectrice $[CI)$ coupe le côté $[AB]$ en $E$. Montrer que le symétrique de $D$ par rapport à la droite $(AH)$ est le point $E$. Comparer $AD$ et $AE$, puis $ID$ et $IE$. 
\item Démontrer que $BD=CE$.
\end{enumerate}
\it on considère un triangle $AHB$ rectangle en $H$. 
\begin{enumerate}
\item On suppose que $AB=2HB$. Démontrer que l'angle aigu en $B$ est le double de l'angle aigu en $A$ (prolonger $[BH]$ de $HC=BH$).
\item On suppose que par hypothèse l'angle en $B$ soit le double de l'angle en $A$. Démontrer que l'on a $AB=2HB$.
\end{enumerate}
\it \begin{enumerate}
\item Construire un quadrilatère convexe $ABCD$ dans lequel les trois angles en $D$, $A$, et $B$ sont égaux à $100^o$ et les côtés $[AB]$ et $[AD]$ égaux à $20$ mm.
\item Démontrer que le triangle $BCD$ est isocèle et comparer $CB$ et $CD$.
\item Que représente la droite $(AC)$ pour le segment $[BD]$ et pour les angles $\widehat{DAB}$ et $\widehat{BCD}$ ?
\end{enumerate}
\it On considère un angle aigu $\widehat{xOy}$ et un point intérieur $A$. \begin{enumerate}
\item Construire les symétriques $B$ et $C$ du point $A$ par rapport à $(Ox)$ et $(Oy)$. Montrer que les points $A$, $B$ et $C$ sont sur un cercle de centre $O$.
\item On mène la droite $(BC)$ qui coupe $(Ox)$ en $M$
et $(Oy)$ en $N$. Comparer les angles $\widehat{OBC}$
et $\widehat{OCB}$, puis montrer que $(AO)$ est la bissectrice de l'angle $\widehat{MAN}$.
\item Soit $D$ le symétrique de $C$ par rapport à $(Ox)$. Démontrer que les trois points $A$, $M$, $D$ sont alignés.
\end{enumerate}



\end{itemize}
 \chapter{Cas d'égalité des triangles}

\begin{itemize}
\it Deux triangles $ABC$ et $A'B'C'$ ont $AB=A'B'$, $\widehat{B}=\widehat{B'}$ et $\widehat{C}=\widehat{C'}$. On mène les hauteurs $(AH)$ et $(A'H')$.\begin{enumerate}
\item Comparer les triangles rectangles $ABH$ et $A'B'H'$. Conséquences ?
\item Comparer les triangles $ACH$ et $A'C'H'$, puis les triangles $ABC$ et $A'B'C'$. 
\item Énoncer le cas d'égalité non classique correspondant.
\end{enumerate}
\it Comparer deux triangles isocèles dans chacun des cas suivants : 
\begin{enumerate}
\item Les bases sont égales et les angles au sommet égaux.
\item Les hauteurs relatives à la base sont égales ainsi que les angles à la base. 
\item Énoncer les cas d'égalité correspondants. 
\end{enumerate}
\it Dans un triangle $ABC$, on prolonge la médiane $(AM)$ d'une longueur $MD=AM$. 
\begin{enumerate}
\item Comparer les triangles $BMD$ et $CMA$, puis évaluer les côtés du triangle $ABD$ par rapport aux côtés $[AB]$, $[AC]$ et à la médiane $(AM)$. 
\item Comparer deux triangles $ABC$ et $A'B'C'$ ayant deux côtés respectivement égaux ainsi que la médiane relative au troisième. Énoncer le cas d'égalité correspondant. 
\end{enumerate}
\it Deux triangles $ABC$ et $A'B'C'$ ont un côté égal $BC=B'C'$ ainsi que les hauteurs $(AH)$ et $(A'H')$ et les médianes $(AM)$ et $(A'M')$. 
\begin{enumerate}
\item Comparer les triangles $AMH$ et $A'M'H'$.
\item On superpose ces deux triangles. En déduire l'égalité des triangles $ABC$ et $A'B'C'$ et énoncer le cas d'égalité correspondant. 
\end{enumerate}
\it Démontrer que dans un triangle isocèle $ABC$ de base $[BC]$. \begin{enumerate}
\item Les médianes relatives aux côtés égaux sont égales.
\item Les bissectrices intérieures des angles $\widehat{B}$ et $\widehat{C}$ sont égales. 
\item Les hauteurs $(BH)$ et $(CK)$ sont égales. Étudier la réciproque. 
\end{enumerate}
\it On prend sur les côtés d'un angle $\widehat{xAy}$ deux points $B$ et $C$ ($AB\neq AC)$. 
La bissectrice de $\widehat{xAy}$ et la médiatrice de $[BC]$ se coupent en $D$. \begin{enumerate}
\item Comparer $DB$ et $CD$ puis les distances $DE$ et $DF$ de $D$ aux côtés de $\widehat{xAy}$.
\item Comparer les triangles $DBE$ et $DCF$, puis les angles $\widehat{BDC}$ et $\widehat{EDF}$.
\item Montrer que $BE$ et $FC$ sont égaux à la demi-différence de $AB$ et $AC$. 
\end{enumerate}
\it Soit un triangle isocèle $ABC$, dans lequel la base $[BC]$ est inférieure aux côtés égaux $[AB]$ et $[AC]$. 
On prolonge $[AB]$ et $[BC]$ de longueurs $BD$ et $CE$ égales à la différence $AB-BC$. \begin{enumerate}
\item Montrer que $BE=AC$. Puis comparer les triangles $ACE$ et $EBD$.
\item Montrer que $\widehat{ADE}=\frac12\left(\widehat{AED}+\widehat{BAC}\right)$.
\end{enumerate}
\it Soient deux points $A$ et $B$ équidistants d'une même droite $(xy)$. On désigne par $M$ et $N$ les pieds des perpendiculaires menées de $A$ et $B$ sur $(xy)$ et par $O$ le milieu de $[MN]$. \begin{enumerate}
\item Comparer les triangles $OAM$ et $OBN$. Conséquences ? 
\item On suppose que $A$ et $B$ soient de part et d'autre de $(xy)$. Montrer que $O$ est aussi le milieu de $[AB]$.
\item On suppose $A$ et $B$ du même côté de $(xy)$. Montrer que la médiatrice de $[AB]$ est la médiatrice de $[MN]$.
\end{enumerate}
\it On considère un triangle isocèle $ABC$ dans lequel 
la médiatrice du côté $[AC]$ coupe le prolongement de la base $[BC]$ en un point $D$. On prolonge $[DA]$ d'une longueur $AE=BD$.\begin{enumerate}
\item Montrer que le triangle $DAC$ est isocèle. Conséquences ?
\item Comparer les triangles $ABD$ et $CAD$. Que peut-on dire du triangle $CDE$ ?
\end{enumerate}
\it La médiatrice du côté $[AB]$ du triangle isocèle $ABC$ coupe la base $[BC]$ (ou son prolongement) en $D$. Le cercle de centre $B$ passant par $D$ recoupe $(AD)$ en $E$. 
\begin{enumerate}
\item Montrer que les triangles $DAB$ et $BDE$ sont isocèles. Comparer les longueurs $AD$ et $BE$ puis les angles $\widehat{ACD}$ et $\widehat{BAE}$ ainsi que les angles $\widehat{ADC}$ et $\widehat{BEA}$.
\item En utilisant le résultat de l'exercice \textbf{25} démontrer l'égalité des triangles $ACD$ et $BEA$ puis que $AE=CD$.
\end{enumerate}
\it \begin{enumerate}
\item Deux triangles $ABC$ et $A'B'C'$ ont $AB=A'B'$, $\widehat{B}+\widehat{B'}= 180^o$ et $\widehat{C}= \widehat{C'}$. Montrer que les hauteurs $AH$ et $A'H'$ sont égales, puis que $AC=A'C'$. 
\item Réciproquement si $AB=A'B'$, $AC=A'C'$ et $\widehat{B}+\widehat{B'}= 180^o$, les angles $\widehat{C}$ et $\widehat{C'}$ sont égaux. (On pourra amener $A'B'C'$ sur $ABC$ de telle sorte que $[A'B']$ coïncide avec $[AB]$ et que $C'$ vienne sur $(CB)$.)
\end{enumerate}
\it Soit un triangle isocèle $ABC$ de base $[BC]$. On
construit une demi-droite $[Bx)$ intérieure à l'angle $\widehat{ABC}$ et une demi-droite $[Cy)$ intérieure à l'angle $\widehat{ACB}$ de telle sorte que les angles $\widehat{ABx}$ et $\widehat{ACy}$ soient égaux. $(Bx)$ et $(Cy)$ se coupent en $M$.
\begin{enumerate}
\item Soient $H$ et $K$ les pieds des perpendiculaires menées de $A$ à $(Bx)$ et à $(Cy)$. Comparer les triangles $ABH$ et $ACK$, puis les segments $[AH]$ et $[AK]$.
\item Établir que la droite $(AM)$ est bissectrice extérieure du triangle $BMC$.
\end{enumerate}
\it Les bissectrices intérieurs des angles en $B$ et $C$ du triangle $ABC$ se coupent en $I$, tandis que leurs bissectrices extérieures se coupent en $J$.\begin{enumerate}
\item Montrer que $I$ et $J$ sont équidistants des trois droites $(AB)$, $(AC)$ et $(BC)$.
\item Démontrer que les trois points $A$, $I$ et $J$ sont alignés.
\end{enumerate}
\it Dans un quadrilatère convexe $ABCD$, on a : $AB=AD$ et $\widehat{B}+\widehat{D}=180^o$. On prolonge $[CB]$ d'une longueur $BE=CD$. \begin{enumerate}
\item Comparer les triangles $ABE$ et $ADC$. Nature du triangle $ACE$. 
\item Démontrer que $(CA)$ est bissectrice intérieure de l'angle en $C$ du quadrilatère.
\end{enumerate}
\it On considère un triangle $ABC$ dans lequel $AB>AC$. La médiatrice de $[BC]$ coupe en $M$ la bissectrice intérieure de l'angle en $A$. Le point $M$ se projette perpendiculairement en $H$ sur $[AB]$ et en $K$ sur le prolongement de $[AC]$.
\begin{enumerate}
\item Comparer les triangles rectangles $MBH$ et $MCK$, puis les triangles rectangles $AMH$ et $AMK$. Démontrer que $AH=AK$ et $BH=CK$. 
\item En déduire que : $AH = \frac12(AB+AC)$ et $BH = \frac12(AB-AC)$.
\end{enumerate}
\it Dans le triangle $ABC$ tel que $AB<AC$, la médiatrice de $[BC]$ coupe en $D$ le côté $[AC]$, en $I$ la bissectrice intérieure de l'angle en $A$ et en $J$ sa bissectrice extérieure.
\begin{enumerate}
\item Montrer que le triangle $DBC$ est isocèle. Que 
représente la droite $(IJ)$ pour l'angle en $D$ du triangle $ADB$ ?
\item Montrer que chacun des points $I$ et $J$ est équidistant des trois droites $(AB)$, $(AD)$ et $(BD)$ puis que $(BI)$ et $(BJ)$ sont bissectrices intérieure
et extérieure de l'angle $\widehat{ABD}$. 
\item Démontrer que les angles $\widehat{ABJ}$ et
$\widehat{ACJ}$ sont égaux tandis que les angles 
$\widehat{ABI}$ et $\widehat{ACI}$ sont supplémentaires.
\end{enumerate}
\end{itemize}









 \chapter{Inégalités dans le triangle}

\begin{itemize}
\it On donne un point $M$ intérieur au triangle $ABC$. 
\begin{enumerate}
\item Montrer que $MA+MB$ est compris entre $AB$ et $CA+CB$.
\item En déduire que $MA+MB+MC$ est compris entre le demi-périmètre et le périmètre du triangle.
\end{enumerate}
\it Soit $O$ le point d'intersection des diagonales d'un quadrilatère convexe $ABCD$. \begin{enumerate}
\item Montrer que chaque diagonale est inférieure au demi-périmètre du quadrilatère.
\item Montrer que $AC+BD$ est supérieure à chacune des sommes $AB+CD$ et $AD+BC$.
\item En déduire que la somme des diagonales est comprise entre le demi-périmètre et le périmètre du quadrilatère.
\end{enumerate}
\it Démontrer que la médiane $AM$ d'un triangle $ABC$ est comprise entre la demi-différence et la demi-somme des côtés $AB$ et $AC$. (On prolongera $[AM]$ d'une longueur égale à elle-même.)
\it Soit un point $A$ et un cercle de centre $O$. Le diamètre passant par $A$ coupe le cercle entre $B$ et $C$. Soit $M$ un point quelconque du cercle. \begin{enumerate}
\item Comparer $AM$ à la somme et à la différence de $OA$ et $OM$.
\item Montrer que $AM$ est compris entre $AB$ et $AC$.
\end{enumerate}
\it Deux points $A$ et $B$ sont d'un même côté d'une droite $(xy)$. \begin{enumerate}
\item Trouver sur cette droite un point $P$, tel que la somme $PA+PB$, soit la plus petite possible (utiliser $A'$ symétrique de $A$ par rapport à $(xy)$.)
\item Que représente $(xy)$ pour l'angle $\widehat{APB}$ ?
\item Montrer que lorsque le point $M$ décrit la droite $(xy)$ la somme $MA+MB$ augmente en même temps que $PM$. 
\end{enumerate}
\it Deux points $A$ et $B$ sont de part et d'autre de $(xy)$. \begin{enumerate}
\item Trouver sur cette droite un point $P$, tel que la différence entre $PA$ et $PB$ soit la plus grande possible (utiliser $A'$ symétrique de $A$ par rapport à $(xy)$). 
\item Que représente $(xy)$ pour l'angle $\widehat{APB}$ ?
\end{enumerate}
\it On donne un angle $\widehat{xOy}$ inférieure à $60^o$ et deux points $A$ et $B$ intérieurs à cet angle. Trouver un point $M$ sur $(Ox)$ et un point $N$ sur $(Oy)$ de façon que le périmètre de la ligne brisée $AMNB$ soit le plus petit possible (utiliser les symétriques $A'$ de $A$ par rapport à $(Ox)$ et $B'$ de $B$ par rapport à $(Oy)$).
\it On considère un triangle $ABC$ et la bissectrice extérieure de l'angle en $A$. \begin{enumerate}
\item Montrer que le symétrique $C'$ de $C$ par rapport à cette bissectrice se trouve sur $(BA)$ et que $AC'=AC$. 
\item Démontrer que pour tout point $M$ de la bissectrice on a : \[ MB+MC>AB+AC.\]
\end{enumerate}
\it On considère un triangle $ABC$ et la bissectrice intérieure de l'angle en $A$. \begin{enumerate}
\item Montrer que le symétrique $C'$ de $C$ par rapport à cette bissectrice se trouve sur $(AB)$ et que $AC'=AC$. 
\item Démontrer que pour tout point $M$ de la bissectrice on a : \[ MB-MC>AB-AC.\]
\end{enumerate}
\it Dans un triangle $ABC$ on prolonge la médiane $[CM]$ d'une longueur $MD=CM$. \begin{enumerate}
\item Comparer les triangles $MBC$ et $MAD$ et montrer que $\widehat{BAD}=\widehat{ABC}$.
\item Soit $\widehat{BAx}$ l'angle extérieur en $A$ au
triangle $ABC$. Démontrer que $\widehat{BAx}>\widehat{ABC}$ et en déduire le théorème : \\
\emph{Dans tout triangle, un angle extérieur est supérieur à tout angle intérieur non adjacent.}
\end{enumerate} 
\it On considère deux triangles $ABC$ et $A'B'C'$ tels 
que $BC=B'C'$ et $\widehat{B}=\widehat{B'}$. On superpose ces deux triangles en amenant $B'$ en $B$, $C'$ en $C$ et $A'$ sur la demi-droite $[BA)$.\begin{enumerate}
\item En utilisant le théorème établi à l'exercice précédent, démontrer que l'une des inégalités $AB>A'B'$, $\widehat{C}>\widehat{C'}$ et $\widehat{A}<\widehat{A'}$ entraîne les deux autres.
\item Qu'arrive-t-il si $\widehat{A}=\widehat{A'}$ ? En déduire le cas d'égalité non classique : \\
\emph{Lorsque deux triangles ont deux angles homologues respectviement égaux et le côté opposé à l'un d'eux égal, ils sont égaux.}
\end{enumerate}
\it Soit un triangle $ABC$ dans lequel $AB<AC$. On prolonge la médiane $[AM]$ d'une longueur $MD=AM$.
\begin{enumerate}
\item Comparer les triangles $MCA$ et $MDB$. Conséquences ?
\item En déduire que : $\frac12(AC-AB)<AM<\frac12(AB+AC)$. Comparer la somme des médianes au périmètre du triangle.
\item Démontrer que l'angle $\widehat{MAC}$ est inférieur à l'angle $\widehat{MAB}$ et que la bissectrice intérieure de l'angle $\widehat{BAC}$ est située à l'intérieur de l'angle $\widehat{MAB}$.
\end{enumerate}
\end{itemize}

 \chapter{Perpendiculaires et obliques}

\begin{itemize}
\it On considère un quadrilatère convexe $ABCD$ dans lequel les angles en $C$ et $D$ sont droits : 
\begin{enumerate}
\item Comparer $AD$ à $AC$ puis $AC$ à $AB+BC$.
\item En déduire que $AD<AB+BC$.
\end{enumerate}
\it Soit un triangle $ABC$ rectangle en $A$. La bissectrice intérieure de l'angle en $B$ coupe $(AC)$ en $D$, et on mène $(DE)$ perpendiculaire en $E$ à $(BC)$.
\begin{enumerate}
\item Comparer $AD$ et $DE$ puis $DE$ et $DC$. 
\item En déduire l'inégalité $AD<DC$. 
\end{enumerate}
\it On considère un triangle $ABC$ dans lequel on mène la hauteur $[AH]$.\begin{enumerate}
\item Comparer $AH$ à $AB$ puis à $AC$ et montrer que $2AH<AB+AC$.
\item En déduire que la somme des hauteurs est inférieure au périmètre du triangle $ABC$.
\end{enumerate}
\it Soit un triangle $ABC$ dans lequel on a $AB<AC$. Soit $[AM]$ la médiane issue de $A$. 
\begin{enumerate}
\item Démontrer que $\widehat{AMB}<\widehat{AMC}$ et que $\widehat{AMB}<90^o<\widehat{AMC}$.
\item En déduire que la hauteur $(AH)$ est du même côté que $(AB)$ par rapport à $(AM)$.
\end{enumerate}
\it On mène la bissectrice $[Oz)$ de l'angle $\widehat{xOy}$ et on porte deux longueurs égales 
$OA$ sur $[Ox)$ et $OB$ sur $[Oy)$.
\begin{enumerate}
\item Que représente $(Oz)$ pour le segment $[AB]$ ?
\item Montrer que pour tout point $M$ intérieur à l'angle $\widehat{xOz}$ on a $MA<MB$.
\end{enumerate}
\it Deux triangles isocèles $OAB$ et $OA'B'$ sont tels que $OA=OB=OA'=OB'$. Soient $H$ et $H'$ les milieux des bases $[AB]$ et $[A'B']$.
\begin{enumerate}
\item On suppose $\widehat{AOB}<\widehat{A'OB'}$. comparer $AB$ et $A'B'$ puis $AH$ et $A'H'$.
\item On suppose $AH<A'H'$. Comparer les angles $\widehat{AOB}$ et $\widehat{A'O'B'}$.
\item En déduire que : \emph{Si deux triangles rectangles ont même hypoténuse, les côtés de l'angle droit sont dans le même ordre de grandeur que les angles opposés.}
\end{enumerate}
\it Deux triangles rectangles $ABC$ et $A'BC$ sont situés du même côté de leur hypoténuse commune $[BC]$. 
On suppose $\widehat{BCA}<\widehat{BCA'}$ si bien que la droite $(CA)$ coupe le côté $[BA']$ en un point $D$ situé entre $B$ et $A'$.\begin{enumerate}
\item Montrer que l'angle $\widehat{CDB}$ est obtus, extérieur au triangle rectangle $\widehat{BAD}$ et que $D$ se trouve également entre $A$ et $C$.
\item En déduire que $BA<BA'$, $CA>CA'$ et $\widehat{ABC}>\widehat{A'BC}$. Étudier les réciproques.
\end{enumerate}
\it D'un point $A$, on mène la perpendiculaire $(AH)$ et les obliques $(AB)$ et $(AC)$ à la droite $(BC)$. \begin{enumerate}
\item Démontrer que l'une des égalités $HB=HC$, $AB=AC$, $\widehat{HAB}=\widehat{HAC}$, et $\widehat{HBA}=\widehat{HCA}$ entraîne les trois autres.
\item En supposant $B$ et $C$ d'un même côté de $H$ et en utilisant le théorème établi à l'exercice \textbf{50}, démontrer que l'une des inégalités $HB<HC$, $AB<AC$ , $\widehat{HAB}<\widehat{HAC}$ et $\widehat{HBA}>\widehat{HCA}$ entraîne les trois autres.
\end{enumerate}
\it On considère un angle $\widehat{xOy}$ et sa bissectrice $[Oz)$. Soit $M$ un point quelconque intérieur à l'angle $\widehat{xOz}$, $N$ et $P$ ses symétriques par rapport à $(Ox)$ et $(Oy)$, $A$ et $B$ ses projections sur $(Ox)$ et $(Oy)$. 
\begin{enumerate}
\item Comparer les triangles isocèles $MON$ et $MOP$ et en déduire que $MA<MB$.
\item Étudier réciproquement la région où se trouve le point $M$ intérieur à l'angle $\widehat{xOy}$ si on a par hypothèse $MA<MB$.
\end{enumerate}
\it On considère deux droites $(AA')$ et $(BB')$ se coupant en $O$ et les bissectrices $[Ox)$ et $[Oy)$ des angles $\widehat{AOB}$ et $\widehat{AOB'}$, ainsi que $[Ox')$ et $[Oy')$ des angles $\widehat{A'OB'}$ et $\widehat{A'OB}$.
\\ Démontrer en utilisant les résultats des exercices \textbf{58} ou \textbf{61} que : 
\begin{enumerate}
\item Tout point $M$ situé à l'intérieur de l'un des angles droits $\widehat{xOy}$ ou $\widehat{x'Oy'}$ est plus près de la droite $(AA')$ que de la droite $(BB')$.
\item Étudier et énoncer la réciproque de cette propriété.
\end{enumerate}
\end{itemize}
 \chapter{Droites parallèles}

\begin{itemize}
\it Démontrer qu'une perpendiculaire $(Ax)$ et une oblique $(By)$ à une même droite $(AB)$ sont sécantes.
\it Démontrer que si deux droites $(\Delta)$ et $(\Delta')$ sont respectivement perpendiculaires à deux droites parallèles $(D)$ et $(D')$, elles sont parallèles entre elles. 
\it Démontrer que, si deux droites $(Ax)$ et $(By)$ sont respectivement perpendiculaires aux côtés d'un angle saillant $\widehat{AOB}$, elles sont concourantes. 
\it Soit un triangle isocèle $ABC$ de base $[BC]$. \begin{enumerate}
\item Montre rque la bissectrice extérieure de l'angle $\widehat{A}$ est parallèle à la base $[BC]$.
\item La réciproque est-elle vraie ? La démontrer.
\end{enumerate}
\it Soient $[AB]$ et $[CD]$ deux diamètres d'un cercle de centre $O$.
\begin{enumerate}
\item Comparer les directions de $[AC]$ et de $[BD]$ avec celle de la bissectrice de l'angle $\widehat{AOC}$.
\item Que peut-on en conclure pour $(AC)$ et $(BD)$ ?
\end{enumerate}
\it Soit un angle $\widehat{xOy}$. Du point $O$ comme centre, on trace un premier cercle qui coupe $[Ox)$ en $A$ et $[Oy)$ en $B$, puis un second cercle qui coupe $[Ox)$ en $C$ et $[Oy)$ en $D$. 
\begin{enumerate}
\item Comparer les directions de $[AB]$ et de $[CD]$ avec celle de la bissectrice de l'angle $\widehat{xOy}$.
\item En déduire que $(AB)$ et $(CD)$ sont parallèles.
\end{enumerate}
\it On considère sur un cercle de centre $O$ qutre points $A$, $B$, $C$ et $D$ dans cet ordre et tels que les arcs $\arc{AB}$ et $\arc{CD}$ soient égaux. 
\begin{enumerate}
\item Montrer que les angles $\widehat{AOD}$ et $\widehat{BOC}$ ont même bissectrice.
\item Comparer les directions de $(AD)$ et de $(BC)$.
\end{enumerate}
\it D'un point intérieur $M$ on mène les perpendiculaires $(MA)$ et $(MB)$ aux côtés de l'angle
droit $\widehat{xOy}$.
\begin{enumerate}
\item Montre que $(MA)$ et $(Oy)$ sont parallèles, et qu'il en est de même de $(MB)$ et $(Ox)$.
\item Quelle est la valeur de l'angle $\widehat{AMB}$ ?
\end{enumerate}
\it On considère un point $M$ pris sur un demi-cercle de diamètre $[AB]$ et de centre $O$. On mène les perpendiculaires $(OH)$ à $(MA)$ et $(OK)$ à $(MB)$.
\begin{enumerate}
\item Montrer que l'angle $\widehat{HOK}$ est droit. Comparer les directions de $(OH)$ et de $(MB)$ ainsi que celles de $(OK)$ et de $(MA)$.
\item Évaluer l'angle $\widehat{AMB}$ et montrer que les angles $\widehat{MAB}$ et $\widehat{MBA}$ sont complémentaires.
\end{enumerate}
\it Soit un quadrilatère $ABCD$ dans lequel les diagonales $[AC]$ et $[BD]$ ont même milieu $O$. On mène les perpendiculaires $(OH)$ à $(AB)$ et $(OK)$ à
$(CD)$. \begin{enumerate}
\item Comparer les triangles $OAB$ et $OCD$, puis les triangles $OAH$ et $OCK$. Conséquences ?
\item Comment sont disposés les points $O$, $H$ et $K$ ? En déduire que $(AB)$ et $(CD)$ sont parallèles. En est-il de même pour $(AD)$ et $(BC)$ ?
\end{enumerate}
\it Étant donné une droite $(xy)$ et un point extérieur $O$, on mène la perpendiculaire $(OH)$ et une oblique $(OM)$ à la droite $(xy)$ puis on construit $H'$ et $M'$ symétriques de $H$ et $M$ par rapport à $O$. \begin{enumerate}
\item Comparer les triangles $OHM$ et $OH'M'$. Conséquences ?
\item Que peut-on dire des directions $(HM)$ et $(H'M')$ ? En déduire que : \\
\emph{Deux droites symétriques l'une de l'autre par rapport à un point sont parallèles.}
\item En déduire une construction de la parallèle menée par $A$ à la droite $(xy)$.
\end{enumerate}
\it On considère deux parallèles $(D)$ et $(D')$ perpendiculaires en $A$ et $A'$ au segment $[AA']$ dont on désigne par $O$ le milieu.\begin{enumerate}
\item Démontrer que la médiatrice $(xy)$ du segment $[AA']$ est un axe de symétrie de la bande $(D, D')$.
\item Établir que $(xy)$ est médiatrice de tout segment $[HH']$ perpendiculaire à $(D)$ et $(D')$.
\item En déduire que tout point $I$ de $(xy)$ est un centre de symétrie des deux droites $(D)$ et $(D')$ et qu'il est équidistant de ces deux droites.
\end{enumerate}
\end{itemize}
 \chapter{Propriétés angulaires des parallèles}


\begin{itemize}
\it Démonter que pour que deux droites soient sécantes, il suffit : 
\begin{enumerate}
\item Que l'une soit perpendiculaire et l'autre oblique par rapport à une troisième.
\item Qu'elles soient perpendiculaires aux côtés d'un angle saillant.
\end{enumerate}
\it Étant données deux parallèles coupées par une sécante, montrer que : 
\begin{enumerate}
\item Les bissectrices de deux angles alternes-internes ou correspondants sont parallèles. 
\item Les bissectrices de deux angles intérieurs d'un même côté sont perpendiculaires. 
\item Énoncer les réciproques de ces propriétés et les démontrer. 
\end{enumerate}
\it Sur les côtés d'un angle de sommet $O$ on porte deux longueurs égales $OA=OB$. Puis extérieurement à cet angle, on construit deux angles égaux $\widehat{OAx}$ et $\widehat{OBy}$ et on porte sur $[Ax)$ et $[By)$ deux longueurs égales $AC=BD$.\begin{enumerate}
\item Comparer les triangles $OAC$ et $OBD$.
\item Montrer que les angles $\widehat{AOB}$ et $\widehat{COD}$ ont même bissectrice.
\item Comparer les directions de $(AB)$ et $(CD)$.
\end{enumerate}
\it On considère deux angles adjacents supplémentaires $\widehat{AOB}$ et $\widehat{BOC}$ et leurs bissectrices $[Ox)$ et $[Oy)$. Par le point $B$, on mène la parallèle à $(AC)$ qui coupe $[Ox)$ en $M$ et $[Oy)$ en $N$.\begin{enumerate}
\item Comparer $MB$ et $OB$ ; puis $NB$ et $OB$.
\item Que représente le point $B$ pour le segment $[MN]$ ?
\end{enumerate}
\it Par le point de rencontre $I$ des bissectrices intérieures des angles en $B$ et $C$ du triangle $ABC$, on mène la parallèle à $(BC)$, qui coupe $(AB)$ en $D$ et $(AC)$ en $E$. \begin{enumerate}
\item Comparer $DB$ et $DI$ puis $EC$ et $EI$.
\item Montrer que $DE=BD+CE$. 
\end{enumerate}
\it Soit un triangle $ABC$ rectangle en $A$. Sur la perpendiculaire en $C$ à $(AC)$ on porte des segments $[CD]$ et $[CE]$ égaux à $[BC]$.\begin{enumerate}
\item Comparer les directions de $(AB)$ et $(DE)$.
\item Que représentent $(BD)$ et $(BE)$ pour l'angle en $B$ ? 
\end{enumerate}
\it Soit un triangle $ABC$. On mène les bissectrices intérieures des angles en $B$ et $C$ qui coupent en $D$ et $E$ la parallèle menée par $A$ à $(BC)$. \begin{enumerate}
\item Comparer $AD$ et $AB$ puis $AE$ et $AC$.
\item Montrer que $DE=AB+AC$.
\item Reprendre le problème avec les bissectrices extérieures.
\end{enumerate}
\it Soit un triangle $ABC$. On mène par le milieu $D$ de $[BC]$ la perpendiculaire à la bissectrice intérieure de l'angle en $A$. Cette perpendiculaire coupe $(AB)$ en $E$ et $(AC)$ en $F$. $(EF)$ coupe alors la parallèle menée par $B$ à $(AC)$ en $G$. 
\begin{enumerate}
\item Montrer que les triangles $AEF$ et $BEG$ sont isocèles.
\item Comparer les triangles $DBG$ et $DCF$. 
\item Démontrer l'égalité des longueurs $BE$ et $CF$.
\end{enumerate}



\end{itemize}
 \chapter{Somme des angles d'un triangle}

\begin{itemize}
\it Soient deux triangles $ABC$ et $A'B'C'$ ayant leurs côtés respectivement parallèles. $(B'C')$ coupe les droites $(AB)$ et $(AC)$ en $D$ et $E$. 
\begin{enumerate}
\item Comparer les angles des triangles $ABC$ et $ADE$, puis ceux des triangles $ADE$ et $A'B'C'$. 
\item Énoncer la propriété qui en résulte pour deux triangles qui ont leurs côtés parallèles.
\end{enumerate}
\it Soient deux demi-droites $[Ax)$ et $[By)$ situées d'un même côté de la droite $(AB)$ et $[Ax')$ et $[By')$ les demi-droites opposées. On suppose $\widehat{BAx}+\widehat{ABy}< 180^o$.
\begin{enumerate}
\item Montrer que les droites $(xx')$ et $(yy')$ sont sécantes et que les demi-droites $[Ax')$ et $[By')$ n'ont pas de point commun.
\item En déduire que $[Ax)$ et $[By)$ se coupent en un point $C$ et énoncer la condition pour que deux demi-droites aient un point commun. 
\end{enumerate}
\it Soit un triangle $ABC$, dont on note les mesures des angles $\widehat{A}$, $\widehat{B}$ et $\widehat{C}$.\begin{enumerate}
\item Évaluer en fonction de $\widehat{B}$ et $\widehat{C}$ l'angle de la hauteur $(AA')$ issue de $A$ et de la  bissectrice intérieure $(AD)$ de l'angle au sommet $A$.
\item Évaluer en fonction de $A$, l'angle $\widehat{BIC}$ des bissectrices intérieures aux sommets $B$ et $C$. 
\item Évaluer en fonction de $\widehat{A}$, l'angle $\widehat{BHC}$ des hauteurs issues de $B$ et $C$. \\
Application numérique : $\widehat{B} = 57^o$, $\widehat{C} = 75^o$.
\end{enumerate}
\it Soit un triangle $ABC$ dans lequel $AC>AB$. On mène la bissectrice de l'angle en $A$ qui coupe $[BC]$ en $D$, puis la perpendiculaire $(BE)$ à $(AD)$. Évaluer en fonction des angles en $B$ et $C$ du triangle : 
\begin{enumerate}
\item Les angles $\widehat{ADB}$ et $\widehat{ADC}$. 
\item Les angles $\widehat{ABE}$ et $\widehat{EBD}$.\\
Application numérique : $\widehat{B}= 68^o$ et $\widehat{C}= 54^o$.
\end{enumerate}
\it Dans un triangle $ABC$, la bissectrice de l'angle $B$ coupe en $I$ la hauteur issue de $A$ et en $D$ la perpendiculaire en $A$ à $(AB)$. 
\begin{enumerate}
\item Évaluer en fonction de $\widehat{B}$ les angles $\widehat{IAB}$, $\widehat{AID}$ et $\widehat{ADI}$.
\item Comparer les segments $[AI]$ et $[AD]$.
\end{enumerate}
\it On considère un cercle $O$ et un diamètre $(xy)$ de ce cercle. D'un point $M$ de ce cercle tel que $\widehat{yOM}<45^o$, on mène $(MH)$ perpendiculaire à $(xy)$ et on construit le point $A$ de $(xy)$ tel que $HA=OH$. La droite $(AM)$ recoupe le cercle en $B$. \begin{enumerate}
\item Montrer que les triangles $MOA$ et $OMB$ sont isocèles. 
\item Calculer les angles $\widehat{OBM}$ et $\widehat{BOx}$ en fonction de l'angle $\widehat{OAB}$.
\end{enumerate}
\it Soit un triangle $ABC$ rectangle en $A$. On prolonge $[CB]$ d'une longueur $BD=BA$ et sur la perpendiculaire en $C$ à $(BC)$, on porte du côté de $A$ une longueur $CE=CA$. 
\begin{enumerate}
\item Calculer en fonction de $\widehat{B}$ les angles $\widehat{BAD}$, $\widehat{ACE}$ et $\widehat{CAE}$. 
\item Évaluer l'angle $\widehat{DAE}$ et montrer que $D$, $A$ et $E$ sont alignés.
\end{enumerate}
\it Soit un triangle rectangle isocèle $ABC$. Par le sommet $A$ de l'angle droit, on mène, extérieurement au triangle une droite $(xy)$, puis les perpendiculaires $(BM)$ et $(CN)$ à $(xy)$ ainsi que la hauteur $(AH)$ du triangle $ABC$. \begin{enumerate}
\item Montrer que $HA=HB=HC$.
\item Comparer les triangles $AMBC$ et $CNA$. Que représente $MN$ pour $BM$ et $CN$ ? 
\item Comparer les triangles $HBM$ et $HAN$ et démontrer que le triangle $HMN$ est rectangle isocèle.
\end{enumerate}
\it On mène la hauteur $(AH)$ issue du sommet de l'angle droit d'un triangle rectangle $ABC$ puis les bissectrices intérieures des angles $\widehat{BAH}$ et $\widehat{CAH}$ qui coupent l'hypoténuse en $D$ et $E$. \begin{enumerate}
\item Évaluer la valeur de l'angle $\widehat{DAE}$.
\item Monter que les triangles $BAE$ et $CAD$ sont isocèles.
\item Comparer $DE$ à la longueur $AB+AC-BC$.
\end{enumerate}
\it Soit un triangle $ABC$. On prolonge $[BC]$ de deux longueurs $BD=BA$ et $CE=CA$. 
\begin{enumerate}
\item Que représente $DE$ pour le triangle $ABC$ ? 
\item Calculer les angles $\widehat{D}$ et $\widehat{E}$ du triangle $ADE$ en fonction des angles $\widehat{B}$ et $\widehat{C}$ de $ABC$.
\item Comparer deux triangles ayant leurs angles respectivement égaux et même périmètre. \end{enumerate}
\it Dans un triangle $ABC$, l'angle aigu $\widehat{B}$ est le double de l'angle $\widehat{C}$. La médiatrice de $[AC]$ coupe $(BC)$ en $D$. 
\begin{enumerate}
\item Montrer que $(AD)$ partage $ABC$ en deux triangles isocèles.
\item Comparer l'angle extérieur au sommet $A$ à l'angle au sommet $C$ dans le triangle $ABC$.
\end{enumerate}
\it Dans un triangle $ABC$, l'angle $\widehat{B}$ est le triple de l'angle $\widehat{C}$. La médiatrice de $[BC]$ coupe $(CA)$ en $D$. 
\begin{enumerate}
\item Montrer que $(BD)$ partage $ABC$ en deux triangles isocèles.
\item Comparer l'angle extérieur au sommet $A$ à l'angle au sommet $C$ dans le triangle $ABC$.
\end{enumerate}
\end{itemize}

% \input{30Parallelogramme}
% \input{31Rectangle-losange-carre}
% \input{32Trapeze-paralleles_equidistantes}
% \input{33Droites_concourantes_du_triangle}
% \input{34Cercle}
% \input{35Position_de_deux_cercles}
% \input{36Cordes}
% \input{37Angle_inscrit}
% \input{38Quadrilatere_inscriptible}
% \input{39Constructions_geometriques}
% \input{40Constructions_de_triangles}
% \input{41Tangentes_et_cercles_tangents}
% \input{42_Polygones_reguliers}
% \input{43Problemes}



%\part{Travaux pratiques et astronomiques}
% \input{44Travaux_pratiques}
% \input{45Coordonnees_equatoriales}
% \input{46Le_systeme_solaire}
% \input{47Eclipses}
 
 	\end{document}
